\section{The State-Basis Discipline}
\label{sec:basis}
\label{sec:state-basis}

The State-Basis Discipline stops a valuation from combining a quantity and a price that were
measured in different corporate-action states of the same instrument. In this framework the
ledger, not the market-data feed, decides which state a unit is in: it reads that state from
the unit's own logged history, and it treats a price from a feed as an outside reading to be
stamped and then either used or refused.

Every corporate action --- a split, a dividend, an index change --- rewrites the frame in which
a unit's numbers mean anything, and the same number can be right in one state and wrong in another. The rule
the discipline enforces is narrow: a quantity and a price may be multiplied only when they are
in the same frame, or after one has been converted into the other's frame by the exact rule the
action itself published. The discipline does not keep the ledger ``in sync'' with the feed, and
it never lets the feed decide a unit's state; where two numbers are in different frames and no
declared conversion connects them, the combination is refused, not guessed.

Here is the error it makes impossible. A wallet holds 1{,}000 shares at \EUR{100}, so its marked
value is \EUR{100{,}000}. The issuer runs a 2-for-1 split: the holding becomes 2{,}000 shares,
and the price reference halves to \EUR{50}. The value is unchanged, because $2{,}000\times50 =
\EUR{100{,}000}$. Now suppose a price feed still carries the last pre-split print, \EUR{100}. A
valuation that multiplies the new quantity by the old price computes $2{,}000\times100 =
\EUR{200{,}000}$ --- a phantom \EUR{100{,}000} that never existed. Each factor is correct on its
own: $2{,}000$ is the real share count, and \EUR{100} was a real trade. Nothing inside the ledger
disagrees, so no reconciliation catches it. The one mistake is combining two numbers that belong
to different frames.

The frame a unit's numbers are quoted in is called its \emph{basis}. A basis makes two numbers
comparable; it does not certify what a holding is worth. A split changes the numbers but not the
value; a spin-off changes past numbers precisely because real value has left the unit for the
child that was spun off.

Unit state already lives in three homes (\S\ref{sec:states}), and the discipline asks one thing
of each; it adds no fourth. \emph{ProductTerms}, the immutable first home, holds the one rounding
rule used when an exact price is rounded to a whole number of minor currency units (cents, for euros); because that rule lives in a single
place, converting a number between bases stays exact. \emph{UnitStatus}, the shared second home,
gains one field: the unit's current basis, which the ledger computes from the unit's own logged
corporate-action events. \emph{PositionState}, the third home, holds the balances, and a balance
is read in the current basis; a balance and a price are brought together only inside a single
valuation, whose basis the ledger fixes.

Observation time cannot stand in for the basis. A feed can lag the moment the ledger marks a
unit \emph{ex} --- the point at which its price detaches from a distribution such as a dividend
--- the ex date can differ between venues, and a late-announced action can change the frame of a
past reading after the fact. The basis is therefore stored explicitly, as a field of the second home, and is
never inferred from a clock.

The whole discipline reduces to one rule, stated formally later as Invariant~\ref{inv:basis},
\emph{single-basis consumption} (\S\ref{sec:basis-invariant}), and referred to by that name
throughout: a valuation runs in one basis, and every balance and every price it uses must be in
that basis --- either because it already was, or because the ledger converted it there by a
declared rule. Each price carries a \emph{stamp} that records which basis it was observed in (\S\ref{sec:basis-ingest}); the check covers every unit named on that stamp, not only the units being
valued. The rule is enforced by the types, not checked after the fact: a quantity and a price in
different bases cannot be multiplied in a well-typed program, and at the boundary where types no
longer reach, the combination is refused rather than guessed (\S\ref{sec:basis-types}).

The rest of this section builds the machinery that makes the rule hold: the basis coordinate and
the two checks that keep it honest; the declared conversion operators; the door through which
outside prices enter and are stamped; the typed read seam; and what happens when a basis is
missing. The same rule governs every place a mark is consumed downstream --- portfolio valuation
(\S\ref{sec:valuation}), the futures lifecycle (\S\ref{sec:futures}), and managed accounts
(\S\ref{sec:managed}).

\subsection{The basis coordinate}
\label{sec:basis-coordinate}

A unit's basis changes only at declared events, never by inference. A boundary
event is the logged fact that a unit's frame changed. It records three things:
which unit changed, when the change takes economic effect, and how each kind of
stored datum converts across the change. The rule for each datum kind is a
declared conversion --- the formal conditions and the code also call it an
\emph{arrow}. Two sources declare boundary events, and both enter through the
log: corporate-action transactions on held units, and external basis notices
with no lifecycle leg, such as index divisor maintenance.

\begin{definition}[Boundary event]
\label{def:boundary-event}
A \emph{boundary event} on unit $u$ is a logged event carrying the unit, an effective
time $t_{\mathrm{eff}}$, a \emph{declaration} --- a finite partial map from datum kinds
to operator specifications (\S\ref{sec:basis-operators}) --- and provenance. Boundary
events are content-addressed; \emph{bid} denotes the address. The two sources are
corporate-action lifecycle transactions on held units (\S\ref{sec:contracts}) and
external basis notices with no lifecycle leg --- index divisor maintenance, composition
true-ups --- entering as logged observation events with a canonical handler.
\end{definition}

The 2-for-1 split enters as one such event: it names the unit $u$, carries the
split's effective time, and declares \texttt{Scale}~$2$ --- the quantity factor ---
for each affected \emph{price-like} datum kind, so a stored price of 100 reads as 50
on the far side. The doubling of the share count travels not as a conversion but as
the entitlement moves of the same transaction ($1{,}000\to2{,}000$), welded to the
declaration by the admission check stated later in this subsection
(Principle~\ref{prin:invariance-weld}); a datum kind the declaration does not
mention has no conversion at all and does not cross the boundary. The event's
\emph{bid} is a content address, computed from the event's content alone: the same
event always has the same id, and the id never depends on when the event was booked.

Basis ids name a unit's frames. The \emph{origin} is the frame a unit starts in
at registration. Each boundary event opens a new frame, named by the event's
own id. The \emph{chain} is the unit's boundary events sorted by when they take
effect, not by when they were booked. Two boundaries can take effect on the
same day; for that case the attested notice declares an intra-day precedence
--- a declared rank that says which same-day boundary applies first, required
by the same-$t_{\mathrm{eff}}$ rule W4 (\S\ref{sec:basis-failure}) whenever the
order matters.

\begin{definition}[Basis id, chain, effective order, epoch]
\label{def:basis-id}
A \emph{basis id} of $u$ is either the origin of $u$ or a boundary id
of $u$. The \emph{basis chain} of $u$, relative to a replay view, is the sequence of
$u$'s boundary events known to that view in \emph{effective order}: the lexicographic
total order on
\[
(t_{\mathrm{eff}},\ \mathit{prec},\ \mathit{bid}),
\]
where $\mathit{prec}$ is the intra-day precedence declared in the attested notice
(\S\ref{sec:basis-failure}, W4) and $\mathit{bid}$ is a deterministic final tie-break.
The \emph{epoch ordinal} of a basis id is its position in the chain --- a derived,
per-view quantity, never stored.
\end{definition}

After the split, $u$ has two basis ids: the origin $b$ and the split's id $b'$,
and the chain holds one event. A back-dated boundary discovered later --- one
whose effective time precedes the split --- inserts \emph{before} the split in
the chain. The insertion shifts $b'$ one position later, so its epoch ordinal
changes; the id $b'$ itself does not. This is why a stamp is a basis stamp, not
a timestamp, and why stamps and state store basis ids, never ordinals: ordinals
renumber under retro-effective insertion, while content addresses are stable.

The effective order is total on every log that can arise, so replay is deterministic.
Ordering keys on effective time first, then on the precedence the two notices declare;
only when both of these coincide does a content id $\mathit{bid}$ break the remaining
tie. That final tie-break would change the numbers in just one situation: two boundaries
on the same unit, at the same effective time and with no declared precedence between
them, whose conversions do not commute, so that applying them in the two possible orders
would give different results. The ledger does not let the $\mathit{bid}$ silently decide
that case.
The W4 rule requires the two notices to declare which comes first; if neither does, the
colliding notice is refused and the unit is held in \emph{pending-transition} --- no
conversion crosses it --- until the conflict is resolved. So the order in which a
boundary was booked never affects the result; only its declared effective position
does.

The coordinate lives where every other unit-level status fact lives: as one field of
\texttt{UnitStatus}, written through the existing status machinery.

\begin{definition}[The basis field]
\label{def:usbasis}
A unit's basis is one field of its status, \texttt{usBasis :: BasisId}, declared with
\texttt{UnitStatus} in \S\ref{sec:states-construction}. The field is total from
registration onward: \texttt{defaultStatus}\,$u$ carries \texttt{Origin}\,$u$
(\hyperref[cond:C5]{C5}), so every registered unit always has a basis. Only one
operation writes it --- the \texttt{SetBasis} case of the closed \texttt{StatusWrite}
set --- and \texttt{applyStatus} applies that write as an outright replacement, never an
increment. So the field is last-write-wins and idempotent (P6), exactly like its sibling
status fields. A separate settle mark, \texttt{usLastSettle}, records the basis a unit
was last settled at; it is dimensionally a \texttt{Price}.
\end{definition}

At registration $u$'s status carries \texttt{usBasis} $= b$, the origin; committing
the split writes \texttt{SetBasis}\,$b'$, and every later read returns $b'$. The
origin is fixed at registration; no later boundary rewrites it. Re-applying
\texttt{SetBasis}\,$b'$ to a status already at $b'$ is the identity, so P6 is
untouched and no counter is stored.

The intended invariant is simple: \texttt{usBasis} equals the effective-order tip
of the chain as known to the view. The field does not maintain this invariant by
itself. Replay folds the log in booking order (\S\ref{sec:lifecycle}), while the
chain is projected in effective order, and the two orders disagree exactly on
retro-effective insertions. Under naked last-write-wins, a back-dated boundary
booked after the split would overwrite $b'$ with a mid-chain basis --- the
coordinate would move backward. The invariant is therefore welded at the door,
not assumed.

\begin{principle}[Tip weld]
\label{prin:tip-weld}
\texttt{applyTx} admits a transaction carrying \texttt{SetBasis}\,$b$ on unit $u$ only if
$b$ is the last element, in the effective order of
Definition~\ref{def:basis-id}, of $u$'s chain after insertion of the boundary the
transaction itself carries. The condition is checkable at admission from the committed
log and the transaction alone.
\end{principle}

For a chain-tip boundary, the common case, the mandated write is the boundary's
own id: the split writes $b'$. For a retro-effective boundary, the
post-insertion tip is the already-committed tip. The mandated write therefore
re-asserts the standing tip, still $b'$; this is idempotent under P6, and the
coordinate never moves backward: \emph{basis regressed by a late notice} is an
unrepresentable ledger state. The inserted boundary changes the mid-chain
conversions, but conversions are per-view projections of the log, never stored
(Definition~\ref{def:basis-category}), so every subsequent snapshot recomputes
them. Observations stamped in the window repair through re-stamp events
(\S\ref{sec:basis-time-travel}); moves repair through the existing compensating
path. The booking-order status fold and the effective-order chain projection
therefore agree on \texttt{usBasis} at every log prefix.

A corporate action acts on two planes at once, and the machinery keeps them
apart. The split's quantity leg travels as moves while its price rule travels
as the declaration; that separation is general.

\begin{principle}[Two planes]
\label{prin:two-planes}
Entitlements and elections are per-holder consequences, and they settle as
moves: this is the \emph{entitlement plane}. The \emph{quotation plane}
records what the market decides about the quoted series: each boundary
carries exactly one declaration, read identically by every holder, however
elaborate the election. Operators live on the quotation plane. A declaration
never depends on any holder's election, and an election never produces a
conversion.
\end{principle}

The second weld ties the coordinate to the economics. A basis change and the
entitlement moves that make it value-neutral travel in one transaction, and the
transaction is admitted only if the two halves agree. In the identity below,
$f$ is the quantity factor the declaration carries, and $c$ is the per-share
cash the boundary books.

\begin{principle}[Invariance weld]
\label{prin:invariance-weld}
A basis-changing corporate action is one atomic transaction
(\hyperref[cond:C3]{C3}) carrying its entitlement moves, its \texttt{SetBasis},
and its declaration. \texttt{applyTx} admits the transaction only if the
booked moves and the declaration jointly satisfy the invariance identity of
Theorem~\ref{thm:basis-value-invariance}, in the case the notice declares.
The \emph{uniform} case covers every holder-independent declaration --- every
class of the taxonomy of \S\ref{sec:basis-cdm} except the elective class ---
and its check is per holder. Each holder's quantity legs must realise the map
$q\mapsto q\cdot f$: under a split the holder's units scale by $f$ and no
cash moves. Each holder's booked cash legs must equal $q\cdot c$: under an
ordinary dividend the holder's units stand and a cash leg of $q\cdot c$ is
booked. Two refinements to these equalities, cash-in-lieu and withholding,
are stated after this principle. The \emph{elective} case
(\texttt{CAMergerElective}, \S\ref{sec:basis-cdm}) checks the same identity
in aggregate at the confirmed blend; its statement accompanies the worked
elective merger (\S\ref{sec:basis-worked}).
Case selection reads declared notice data only --- the notice's class and
election structure --- never a name or a heuristic. Absent a confirmed
elective declaration, the uniform case applies: the door fails closed into
the strict per-holder weld.
\end{principle}

Two refinements qualify the uniform equalities. Each is an explicit cash leg
of the same transaction, never an adjustment outside it.

\begin{itemize}[nosep]
\item \emph{Cash-in-lieu.} A fractional entitlement is paid as real cash in
place of the fraction, as an explicit leg of the same transaction. The
quantity and cash equalities then hold up to exactly the legs so declared.
\item \emph{Withholding.} Where tax is withheld at source, two cash legs
travel in the same transaction: the holder's net cash leg, and a second leg,
booked on that holder's entitlement, paying the withheld amount to the
designated tax wallet. The two sum to $q\cdot c$. Absent an explicit
withholding leg, the gross identity applies.
\end{itemize}

The split declares $f=2$ and $c=0$: the weld admits it only with quantity legs
taking 1{,}000 shares to 2{,}000 and no cash leg. A \EUR{2} dividend on 1{,}000
shares declares $f=1$ and $c=2$: the weld admits it only with a cash leg of
$1{,}000\times2=\EUR{2{,}000}$. The check needs no market data. It is the same
admission idiom as \texttt{move}'s positivity gate. Consequence: ``positions
doubled but basis not advanced,'' ``basis advanced but declaration missing,''
``declaration inconsistent with moves,'' and, by the tip weld, ``basis
regressed by a late notice'' are all unrepresentable ledger states.

The uniform case is the weld of Principle~\ref{prin:invariance-weld} at full
strength, holder by holder: a split that misallocates shares between holders
while totals match is refused at the door. An elective boundary cannot be
checked holder by holder, because no single holder's legs match the blend ---
that is what an election means. The door therefore checks the quotation-plane
fact (Principle~\ref{prin:two-planes}): the blend identity at unit level,
stated as the elective aggregate weld with the worked elective merger of
\S\ref{sec:basis-worked}. Per-holder allocation is checked separately, as
economic correctness: a wrong allocation is not refused at the door but
caught by recomputing the transaction from the log and comparing. Each
holder's legs are compared
against that holder's logged election and the confirmed proration; elections
and proration are logged, attested data, so the comparison is recomputable
from the log, and differential recomputation (P27) performs it. Each case
keeps its own witness: the per-holder weld witness for uniform boundaries,
and elective aggregate invariance (P33) for elective ones. Both run, and
neither substitutes for the other.

\subsection{The adjustment operators and their composition law}
\label{sec:basis-operators}

Each corporate-action event carries its own adjustment operator as data --- a 2-for-1
split carries \texttt{Scale}~$2$, a \EUR{2} dividend carries \texttt{Shift}~$-2$ ---
and one generic evaluator applies whatever the event declares. The system never
derives an adjustment from an event's name, and never computes one from a model: an
event that declares no operator adjusts nothing --- the safety clause of the
confirmation gate's pending window (\S\ref{sec:basis-cdm}). An operator acts on every declared
field at once, each by the map its dimension fixes
(Definition~\ref{def:dimension-declaration}): under \texttt{Scale}~$f$, quantities
multiply by $f$ and prices divide by $f$; under \texttt{Shift}~$x$, prices move by
$x$; cash and dimensionless values never move. A \EUR{2} dividend lowers the price by 2, so
it is \texttt{Shift}~$-2$: 102 becomes 100, and the matching cash leg pays the 2 to
holders.

\begin{definition}[Operator specification]
\label{def:opspec}
A boundary's declaration is a map that sends each affected datum kind to one operator
from the list below. Every parameter comes from the declared terms of the notice,
never from the event's classification: a 5\% stock dividend that declares $f=21/20$
is a scale, whatever its name. The list is closed: extending it is a change to this
specification, never a change to data, and a declaration naming an operator outside
it is refused at the door.
\end{definition}

\begin{lstlisting}[language=Haskell]
data OpSpec
  = Scale Rational       -- f /= 0;  invertible
  | Shift Rational       -- invertible
  | Subst UnitId Rational -- r /= 0; cross-unit succession, invertible
  | Recompose Sigma      -- invertible iff sigma is declared bijective
  | AId                  -- identity: DECLARED, never assumed
  | Pending BoundaryRef  -- boundary known, parameter unpublished: no arrow yet
  | Terminal             -- series ends: no arrow in either direction
  deriving (Eq, Show)
-- Side conditions (non-zero, bijectivity) are enforced by unexported smart
-- constructors at the single admission door (applyTx), the `move` idiom: the
-- invalid spec is never stored, so the evaluator is total over stored values.
\end{lstlisting}

\noindent
The map is partial on purpose: a datum kind outside the declaration's domain has no
conversion and cannot cross that boundary, so a declaration covering only the spot
kind converts spot and refuses every other kind --- the door fails closed rather than
guess. \texttt{AId} must be declared; it is never a default. All data-plane
arithmetic is exact rational. The single crossing to integer minor units is
\texttt{toPrice}, whose rounding rule is fixed once in \emph{ProductTerms}
(round-half-even, one site). Composition of exact maps is therefore
path-independent, and caching intermediate values is sound.

A \texttt{Recompose}~$\sigma$ whose target names successor units declares one
forward conversion onto a jointly-stamped basket of the successors; the
basket mechanics are stated with the spin-off worked example
(\S\ref{sec:basis-worked}).

A unit's bases are its basis ids: the origin, then one new basis after each boundary.
The declared operators connect them. Each boundary gives, for each datum kind it
declares, one conversion from the basis before it to the basis after it, and carrying
a datum across several boundaries is the composition of those conversions in order.

\begin{definition}[Basis conversions and their composition]
\label{def:basis-category}
The bases are the basis ids of all units. For each datum kind $k$, the declared
conversions are these: for each boundary event $g$ with $k$ in the domain of its
declaration, one conversion from the basis before $g$ to $\mathit{bid}(g)$, acting by
the interpretation of the declared operator; and for each declared succession
(\texttt{Subst}, welded to \texttt{usSupersededBy}, \hyperref[cond:C8]{C8}), one
cross-unit conversion from the predecessor's final basis to the successor's origin.
A reverse conversion exists exactly where the declared operator is invertible;
\texttt{Pending} and \texttt{Terminal} give no conversion. Within one unit the chain
is totally ordered and successions are unique, so between any two bases on one
(possibly succession-extended) chain there is at most one conversion per datum kind.
\end{definition}

Uniqueness of the conversion makes multi-boundary transport a computation, not a
choice: the composite across several boundaries is the composition of the boundary
conversions, in effective order, and it exists exactly where every one of them does.
A \texttt{Recompose} targeting successor units adds one cross-unit basket
step to this law (\S\ref{sec:basis-worked}); each successor's own chain
then composes as above.

\begin{proposition}[Composition law]
\label{prop:composition-law}
For basis ids $b\le b'$ on the (possibly succession-extended) chain, with intervening
boundaries $g_1,\dots,g_n$ in the effective order of Definition~\ref{def:basis-id},
\[
A_k(b\to b') \;=\; \mathrm{interp}(g_n)\circ\cdots\circ\mathrm{interp}(g_1),
\qquad\text{defined iff $k$ is in every $g_i$'s declaration;}
\]
\[
A_k(b'\to b) \;=\; A_k(b\to b')^{-1}
\quad\text{defined iff every intervening conversion is invertible;}
\]
\[
A_k(b\to b)=\mathrm{id},
\qquad
A_k(b\to b'') \;=\; A_k(b'\to b'')\circ A_k(b\to b').
\]
\end{proposition}

The declared operator is written into the log once and consumed twice. The
two consumers share no code path and run at different times: the
corporate-action processor consumes it once, at the effective time; the read
seam consumes it on every later read.

\begin{principle}[One declaration, two consumers]
\label{prin:one-decl-two-uses}
A boundary's declared operator is consumed twice: at the effective time, by
the corporate-action processor that produces the boundary's one-time
transactions; and from then on, at the read seam, as the standing conversion
applied to every pre-boundary arrival. Neither consumer interprets the event;
both apply the same logged declaration.
\end{principle}

The first use is the effective-time work. At $t_{\mathrm{eff}}$ the
corporate-action processor (\S\ref{sec:processor-contract}) reads the
declaration --- for the 2-for-1 split, quantity factor $f=2$ and no cash ---
and produces the transactions the door admits: entitlement moves taking
1{,}000 shares to 2{,}000; a terms append taking the option series' strike
from \EUR{100} to \EUR{50} (\hyperref[cond:C6]{C6}); a terms append taking the
performance option's stored initial spot from \EUR{100} to \EUR{50}. Which
fields move, and by what map, is fixed by the unit's adjustment schedule
(\hyperref[cond:C14]{C14}, Principle~\ref{prin:schedule-authority}). Each
follow-on transaction carries the boundary's id in \texttt{txCause}. That
field is provenance, never control: no admission, valuation, or transport
decision reads its value. It is necessary --- it is what lets an auditor walk
from the \EUR{50} strike back to the notice that caused it --- and it is
insufficient, because it names the cause and not the conversion. Audit needs
the tag; arithmetic needs the operator.

The second use never ends. Prices do not stop arriving because a split
happened. The 09:59 print for a 10:00 split carries a pre-boundary stamp; so
do last month's history pulled next week, a vendor's late file tonight, and
every observation inside a historical view that \texttt{clone\_at($t$)}
rebuilds (\S\ref{sec:basis-time-travel}). The read seam converts each one by the composite of
Proposition~\ref{prop:composition-law}: the declared conversions, in effective
order, across however many boundaries the stamp lies behind. Nothing is
reprocessed and nothing is configured. A backtest also reads the other way,
expressing post-boundary data in an earlier frame, and the reverse conversion
exists exactly where the declared operator is invertible
(Definition~\ref{def:basis-category}): \texttt{Scale} and \texttt{Shift}
always invert; \texttt{Recompose} inverts only where its $\sigma$ is declared
bijective; \texttt{Pending} and \texttt{Terminal} give no conversion in
either direction.

This is why the declaration is an operator and not a name. ``Stock split'' is
an open name: consuming it means interpreting it, identically, in every
consumer, forever. \texttt{Scale}~2 is an entry in a closed menu: consuming it
means applying it. With a bare tag the second use is impossible, because
nothing in the log says divide by 2. The first use is unverifiable, because
the log would show a strike of \EUR{50} and no longer carry the fact that lets
anyone check that \EUR{50} is \EUR{100} halved.

One June-1 timeline carries both uses. The unit's chain already holds May's
\EUR{2} dividend, \texttt{Shift}~$-2$ ($b\to b_1$). At 09:59 a vendor prints
100, stamped $b_1$ under the source's convention on file. At 10:00 the split
takes effect ($b_1\to b_2$): the processor consumes the declaration once, and
the door admits its transaction --- 2{,}000 shares, the \texttt{SetBasis}, the
two terms appends, each tagged with the boundary's id. At 10:03 a valuation
asks for spot; the seam reads the 09:59 print, one boundary behind, and
returns $100/2=50$. Next week a history pull fetches May prints stamped at the
origin $b$: a spot of 102 crosses both boundaries, and the composite is fixed
by effective order, $(102-2)/2=50$. The reversed composite, $102/2-2=49$, is
not a valid conversion: no pair of basis ids has it as its unique conversion,
and no exported constructor can forge it. Had the dividend and the split
shared one $t_{\mathrm{eff}}$, the pair would be inadmissible without a
declared mutual precedence; the 49-versus-50 ambiguity is refused at the door,
never resolved by arrival order. Tonight's late vendor file, rows stamped
$b_1$, converts the same way on read, and so does every later
\texttt{clone\_at($t$)}. On June 1 exactly one component did one-time work:
the processor, once.

The theorem states the economics the invariance weld enforces: crossing a boundary
changes the numbers but not the value, except for exactly the cash the boundary
itself books. In its statement, $f$ is the declared quantity factor and $c$ is the
cash per share the boundary books. The split declares $f=2$, $c=0$: a wallet of
$1{,}000$ shares at 100 becomes $2{,}000$ at 50, and
$2{,}000\times50=\EUR{100{,}000}$ --- value unchanged, no cash booked. The dividend
declares $f=1$, $c=2$ on $1{,}000$ shares priced at 102: the price becomes 100 and
the same transaction books $1{,}000\times2=\EUR{2{,}000}$ of cash ---
\EUR{102{,}000} before, \EUR{100{,}000} in shares and \EUR{2{,}000} in cash after.

\begin{theorem}[Lifecycle value invariance across boundaries]
\label{thm:basis-value-invariance}
Let a boundary declare price action $p\mapsto(p-c)/f$ with quantity legs
$q\mapsto q\cdot f$ and cash legs $q\cdot c$ (up to explicit cash-in-lieu and
withholding legs, each booked in the same transaction), as
Principle~\ref{prin:invariance-weld} requires. Then
\[
(q\cdot f)\cdot\frac{p-c}{f} \;=\; q\cdot p \;-\; q\cdot c :
\]
marked value is continuous through every boundary up to exactly the cash the same atomic
transaction books.
\end{theorem}

\begin{proof}
The identity is algebraic; the weld makes its hypotheses admission conditions, so every
committed boundary satisfies them. For an elective boundary the weld's hypotheses hold
in aggregate, and the identity applies to the aggregate position by linearity:
$(\sum q\cdot f)\cdot\frac{p-c}{f}=\sum q\cdot p-\sum q\cdot c$, with $(f,c)$
holder-independent because operators live on the quotation plane
(Principle~\ref{prin:two-planes}). Fractional entitlements land as explicit cash-in-lieu
moves in the same transaction, so the identity holds exactly, down to the last minor unit: any leftover fraction is
paid as real cash in the same transaction, and is never lost to rounding.
Withholding likewise: the withheld amount is real cash to the tax wallet in
the same transaction, so the holder's net and the withheld part partition
$q\cdot c$ exactly --- recorded, never leaked.
\end{proof}

\noindent
For basis boundaries this proves the lifecycle value-invariance property stated in
restricted form in \S\ref{sec:intro}: the price jump at the boundary exactly offsets
the booked cash. The theorem's scope follows the weld's two cases: per-position
continuity across uniform boundaries, and aggregate continuity across elective
boundaries. Across an elective boundary the per-holder differences are exactly
the elected consideration, and that consideration is recorded in the log. On
the worked merger of \S\ref{sec:basis-worked} the differences are
$\pm\EUR{1{,}000}$ per holder and sum to zero.

The theorem speaks of one quantity, one price, and the booked cash. A
product's terms carry more numbers than that: the vanilla call carries a
strike and a deliverable per contract; the performance option carries a stored
initial spot; a single datum may carry several components, as an implied-vol
point carries a strike coordinate and a vol value. Every one of these must
move, or stay, at every boundary, and which it does is not a property of the
event. It is a property of what the number measures.

\begin{definition}[Dimension declaration]
\label{def:dimension-declaration}
Every terms field and every component of a registered datum kind declares
exactly one dimension at its definition: quantity-of-$u$, price-in-basis,
cash, or dimensionless. A boundary declaring quantity factor $f\neq0$ and
per-share cash $c$ acts on each field by the map its dimension fixes --- the
\emph{per-dimension action} of the table below.
\end{definition}

\begin{center}
\begin{tabular}{@{}llll@{}}
\toprule
\textbf{Dimension} & \textbf{Measures} & \textbf{Action under $(f,c)$} &
\textbf{Examples} \\
\midrule
quantity-of-$u$ & units of $u$ held or delivered & $q\mapsto q\cdot f$ &
position; deliverable per contract \\
price-in-basis & \EUR{} per unit of $u$, in a frame & $p\mapsto(p-c)/f$ &
spot; strike; stored initial spot \\
cash & \EUR{}, in no frame & fixed & booked dividend cash; premium paid \\
dimensionless & a pure number & fixed & moneyness; a vol value; a proportion \\
\bottomrule
\end{tabular}
\end{center}

The dimension is the constructor: the only way to build a terms value is to
build it as a quantity, a price, cash, or a dimensionless number, so
declaring the dimension is the same act that creates the field. An
unclassified field is therefore unconstructible, not merely invalid:

\begin{lstlisting}[language=Haskell]
data Dim        = DimQuantity | DimPrice | DimCash | DimNone
data TermsValue = TVQty Rational | TVPrice Rational
                | TVCash Rational | TVFree String   -- dimension = constructor

transportField :: Rational -> Rational -> TermsValue -> Maybe TermsValue
-- the per-dimension action: quantities x f, prices (p-c)/f, cash and
-- dimensionless fixed; f == 0 is refused before any field is touched.
\end{lstlisting}

\noindent
This extends the scalar discipline of \S\ref{sec:futures}, under which
quantities and cash add and prices do not. Here the same separation, plus the
dimensionless case, becomes a declaration every field must make before any
boundary exists to test it. A field with no dimension cannot be written into
any terms version, so ``the strike was never classified'' is not a latent bug.
It is a value that cannot be built.
The action column is the re-expression map of the declared boundary. Whether
a given corporate-action class applies that map to a given terms field is a
separate, per-product question; the adjustment schedule of
\S\ref{sec:basis-schedule} answers it, through a per-class default reaction
called the \emph{class default}. A quote is not a terms field and faces no
such question: it always crosses by the action column, at the ingest seam
(\S\ref{sec:basis-ingest}).

The vanilla call shows the action. The split declares $f=2$, $c=0$. The strike
is price-in-basis, so \EUR{100} reads \EUR{50}; the deliverable is
quantity-of-$u$, so 100 shares read 200. At any post-split spot $s$ the
position pays $200\cdot\max(s-50,0)$; at the corresponding pre-split spot $2s$
it paid $100\cdot\max(2s-100,0)$ --- the same number, for every $s$. Value
invariance holds field by field, from two declarations and no option-specific
rule.

The performance option shows what the declaration buys. It pays on the ratio
$S_T/S_0$, with the initial spot $S_0=\EUR{100}$ stored as a terms field.
$S_0$ measures euros per share in a frame, so it is declared price-in-basis,
and the split takes it to \EUR{50}; spot observations cross by the same map. A
post-split valuation re-derives the ratio from transported inputs and gets
$(S_T/2)\,/\,(S_0/2)=S_T/S_0$: invariant automatically. No rule was ever
written for performance options. The field declared what it measures, and the
dimension did the rest.
The dimension settles the split; the ordinary dividend falls to the class
default. Under an ordinary dividend the class default
(\S\ref{sec:basis-schedule}) holds $S_0$ fixed. A performance confirmation
fixes $S_0$ at inception, and holding it fixed on an ordinary dividend keeps
exactly that contract term. The ex-date price drop then shows up in the ratio
through fresh prints of $S_T$, never through any change to $S_0$.

\begin{theorem}[Dimension invariance across boundaries]
\label{thm:dimension-invariance}
Let a boundary declare quantity factor $f\neq0$ and per-share cash $c$, as
Principle~\ref{prin:invariance-weld} requires. Under the per-dimension action, in the
same boundary, each price-dimension value maps by $p\mapsto(p-c)/f$, each
quantity-dimension value maps by $q\mapsto q\cdot f$, and cash and
dimensionless values are fixed. Hence for every quantity--price pair a
product's terms and data put together,
\[
(q\cdot f)\cdot\frac{p-c}{f}\;=\;q\cdot p\;-\;q\cdot c:
\]
the identity of Theorem~\ref{thm:basis-value-invariance} holds field by field,
for every product, at every boundary on the closed operator menu.
\end{theorem}

\begin{proof}
The identity is algebraic in $(f,c)$ and never mentions which field carries
$q$ or which carries $p$. The action is fixed by the dimension alone, the
dimensions are a closed four-element set, and every field carries one by
construction. The proof therefore quantifies over four dimensions, once ---
never over products or event classes. A product registered next year adds
fields, not cases.
\end{proof}

\noindent
The theorem moves each value by the dimension it declares; it does not judge
whether the declaration is true. For an observation --- a datum-kind entry ---
the correctness of the declaration is a registration fact. It is fixed at the
point where the kind's dimension is declared, and the per-entry invariance
witness the entry carries checks it there (\S\ref{sec:basis-kinds}). If a
declared dimension is wrong, the theorem still applies that dimension's map,
and the resulting settlement pays the wrong amount. That mispayment is what
the invariance witness, not this theorem, is built to catch.

The declarations relocate the expert work; they do not eliminate it. Someone
who knows the product still decides, field by field, what each number
measures, and that judgement is as hard as it ever was. But it is made once
per field, at definition, as a declaration --- not once per product per event
class, as adjustment code. The relocation changes the count: $N$ products
crossing $M$ event classes need $N$ sets of field declarations and one
class-default table of constant size, declared once in this specification ---
not $N\times M$ hand-written cases. Totality is completeness, not correctness.
A schedule total in the sense of \S\ref{sec:basis-schedule} --- one entry for
every pair of a class and a field --- answers every such pair from its own
declarations, and it answers wrongly wherever a declaration is wrong.
Differential recomputation (P27) checks that
the schedule was applied faithfully; it does not check that the convention
the schedule encodes is the right one. That is why the default table is fixed
in this specification and witnessed, rather than left to each product to
remember.

\paragraph{Pending crystallisation.} Some corporate actions become effective before
their parameter is published: the boundary is known, the number is not. A
\texttt{Pending} boundary advances the basis id at effectiveness. Its moves and
its \texttt{SetBasis} commit in one atomic transaction (\hyperref[cond:C3]{C3}). The
parameter arrives later as a parameter-publication event referencing the boundary id.
That event is a versioned resolution under the \emph{ProductTerms} append discipline
(\hyperref[cond:C6]{C6}). The chain projection then resolves each boundary to its
latest parameter version. No committed element mutates and the epoch does not move,
because a publication is not a boundary. Until publication no conversion crosses that
boundary for the affected kinds, so quarantine is forced (\S\ref{sec:basis-failure}).
The quarantine is correct, because any interim mark is a guess.

Transport applies only to what the ledger observed; what the ledger computed is
recomputed at the target basis, never transported.

\begin{principle}[Re-derivation canon]
\label{prin:rederivation}
Primitive data are transported by declared conversions; derived data are re-derived
from transported inputs.
\end{principle}

The framework never assumes that a derivation commutes with adjustment: whether it
does is a fact about the model that produced the datum, and the model-free directive
--- the rule that the ledger assumes nothing about a pricing model's internals ---
forbids depending on it. The index divisor shows why. An index $I$ has constituents
$A$ at 60 and $B$ at 84, and its published level $(60+84)/D=100$ fixes the divisor
$D=1.44$. $B$ splits 2-for-1, and fresh prices arrive: $A$ at 60, $B$ at 42.
Consuming the stale divisor gives $(60+42)/1.44\approx70.83$, a phantom $-29.2\%$ move built
from individually correct inputs. The ledger refuses it. The divisor is derived, not observed: it was computed from
$A$'s and $B$'s prices, so it means something only while those prices sit in the bases
it was built from. Once $B$ splits, one of those bases has moved and the stored divisor
no longer matches, so it is refused rather than reused. The formal check, which ranges
over every unit the divisor's stamp names, is Invariant~\ref{inv:basis}(ii)
(\S\ref{sec:basis-invariant}).
The repair is a maintenance boundary declaring \texttt{Recompose} with $D'=1.02$,
after which $(60+42)/1.02=100$: the divisor is re-derived, never transported. A
derivation performed inside a snapshot scope consumes transported primitive inputs,
so its result is in that basis as well: the types carry the basis through ordinary
computation, with no separate check (\S\ref{sec:basis-types}).

\subsection{The adjustment schedule}
\label{sec:basis-schedule}

The operators of \S\ref{sec:basis-operators} say how a declared boundary converts
stored data. One question they leave open: for \emph{this} product, what does each
class of event do to each of its terms? A vanilla call and a performance option both
face the same 2-for-1 split, and their strikes react by different rules. That answer
is product data, and it is demanded where all product data is laid down: every
unit's \emph{ProductTerms} carries, from its first version at registration, an
\emph{adjustment schedule} --- for every corporate-action class (\texttt{CAClass},
\S\ref{sec:basis-cdm}) and every terms field, exactly one entry stating how that
field reacts. An entry comes from one of three layers.

\begin{description}[style=nextline]
\item[Class default.] Every terms field already declares its dimension, and this specification
declares once, for every pair of a corporate-action class and a dimension, the default reaction:
quantities scale by $f$; cash and dimensionless fields stand; price-in-basis fields move to
$(p-c)/f$ under every class except \texttt{CADividendOrdinary}, where they
stand.\footnote{In the reference implementation the table is
\texttt{classDefault}; its price-dimension column is
\texttt{priceTermsDefault}.} That one exception follows from the invariance
arithmetic itself, not from market custom; its derivation follows this
list. Quotes are not terms
fields; they cross by the declared operator at the ingest seam
(\S\ref{sec:basis-ingest}), so an old cum print still reads $p-c$ past the
ex-date. A term sheet that deviates declares an override. The table is
written out over the closed \texttt{CAClass}: a new class re-opens it entry
by entry, never through a wildcard.
\item[Declared override.] Where the term sheet deviates from the default, the
deviation is written as a rule in a closed, first-order language: a rule reads the
boundary's declared parameters and the field's current value, combines them with
arithmetic and comparison, and chooses by condition. The language has no recursion
and no partial operation, so every well-formed rule evaluates to exactly one
value. Widening the language is a revision of this specification, never a
property of data, and the totality gate \hyperref[cond:C14]{C14} below is
re-run under the revised language. Example: a term sheet says a special dividend
adjusts the strike only above \EUR{0.50}, and only by the excess. The rule reads the
declared per-share cash $c$ and sets the strike to $K - (c - 0.50)$ when $c > 0.50$,
else leaves it. A \EUR{2} special dividend takes a \EUR{100} strike to \EUR{98.50},
where the class default would give \EUR{98} --- the override exists precisely
because the confirmation says \EUR{98.50}.
\item[Designated discretion.] Where the confirmation empowers a party --- a
calculation agent, an exchange --- to determine the adjusted value, no formula
exists to write down. The schedule then declares \emph{who owes the
determination}, and nothing else. The determination later enters the log as
\emph{resolved terms}: a W4-attested terms append referencing the boundary,
transcribed from the designated party's decision. Resolved terms are the
second admissible source of adjusted terms; the first is the schedule's
evaluator --- the generic component of \S\ref{sec:basis-operators} that
applies declared entries --- and the principle below closes the pair. Until
the determination lands, the affected fields have no adjusted value. The
boundary is \texttt{Pending} for them and consumption quarantines: the
pending-crystallisation discipline of \S\ref{sec:basis-operators}, extended
from a parameter no one has published to a determination no one has yet
made.
\end{description}

The \texttt{CADividendOrdinary} exception derives in three steps. First, the
$-c$ shift is the value-invariance re-expression of
Theorem~\ref{thm:basis-value-invariance}, and it is valid only when the same
transaction that shifts the field also pays this field's holder $q\cdot c$.
Second, an ordinary dividend pays the shareholder, not the option holder, and
an anticipated distribution is already priced into every contractual
reference struck against the underlying --- a strike, a stored initial spot.
Third, shifting such a field on an ordinary dividend would therefore hand its
holder value the boundary has already booked to the shareholder --- an
unearned transfer. So price-in-basis terms fields stand under
\texttt{CADividendOrdinary}. In numbers: a \EUR{100}-strike call crosses an
ordinary \EUR{2} dividend. The wrong reaction re-expresses the strike to
\EUR{98}; the correct one leaves it at \EUR{100}. The \EUR{2} of option value
the wrong shift would create is exactly the \EUR{2} per share the boundary
pays the shareholder --- the same cash booked twice if the strike moved.
Listed and OTC convention --- ordinary distributions adjust no strike;
splits, specials, and successions do --- is thus the arithmetic's own answer,
declared here once for every product.

\begin{principle}[Schedule authority]
\label{prin:schedule-authority}
Adjusted terms enter the log from exactly two sources: the schedule's evaluator
applying the unit's declared entry, and the attested resolved terms of the party the
schedule designates. No other source of adjusted terms exists.
\end{principle}

The schedule is total, and totality is checked where the schedule is written. The
condition below joins the register of \S\ref{sec:states-conditions}:

\begin{description}[style=nextline]
\item[C14 (adjustment-schedule totality).]\hypertarget{cond:C14}{}\label{cond:C14}
Registration, and every terms append, is admitted only if the schedule yields
exactly one well-formed entry for every pair of a corporate-action class and a field
of the version being written; an append that introduces a field extends the schedule
to cover it in the same transaction. The refusal is \texttt{ScheduleIncomplete}.
\end{description}

\noindent
The gate quantifies over the closed \texttt{CAClass}, so it is a finite check.
It is also strict: an ill-formed declared override is refused; it never falls
through to the default. The consequence is that ``the framework does not know how to handle this
event for this product'' is unrepresentable past inception. For every event class
and every field, a registered unit's terms state the answer, the formula that
produces it, or the party who owes it.

Adjustment never edits terms. The evaluator's output is a new terms version
appended under \hyperref[cond:C6]{C6}, carrying the boundary as its cause, so
replay re-runs the evaluator against the logged notices and re-derives every
historical strike.

The schedule governs one plane only: how the unit's single set of terms
reacts. The per-holder entitlement moves of the same transaction are the
other plane, checked by the invariance weld
(Principle~\ref{prin:invariance-weld}).

One duty falls on schedule authors, and it is per class. A listed series
declares the exchange as designated party for exactly the classes on which
the exchange renders adjustment decisions --- special distributions,
spin-offs, elective and prorated events --- and keeps defaults and overrides
for the rest. An ordinary dividend adjusts nothing and the exchange publishes
nothing, so a wholesale designation would hold the series \texttt{Pending} on
every ex-date for a determination no one owes. The schedule is keyed by
(class, field); a designation is an entry, never a blanket.

\subsection{The stamp and the ingest door}
\label{sec:basis-ingest}

A raw vendor number becomes a ledger observation in exactly one way: it passes
through one total ingest function. This is the parse-boundary idiom of
\texttt{move} --- validate at the door, so that everything past the door is
already well formed.

Two terms the door depends on are fixed now. A datum kind is
\emph{registered} when a registration event has declared one dimension for
each of its components (\S\ref{sec:basis-kinds}). An unregistered kind is
unconsumable: the door refuses it outright.

\begin{lstlisting}[language=Haskell]
ingest :: Ledger -> SourceId -> Timestamp -> RawDatum
       -> Either IngestError StampedObs
data StampedObs   -- ABSTRACT: constructor unexported; ingest is the only door.
  -- Carries: value (exact), t_obs, source, and stamp :: Map UnitId BasisId,
  -- inside the attestation envelope (provider key, signature, content address).
\end{lstlisting}

\noindent
The stamp is assigned per (unit, source) pair. A source can lag the unit: after
the 2-for-1 split the unit's chain has advanced to the post-split basis, yet a
source may go on publishing pre-split prices for a while --- the ex-transition
lag of \S\ref{app:pricing-coordination}, whose one-bit quote-ex flag and
hard-coded one-step adjustment are the special case of a chain with a single
boundary. The lag is recorded as a per-source
basis offset, from one of two places: transcribed from
the source's declared convention under signature, or asserted by the gateway
against the committed corporate-action log. The stamp is never defaulted from
the ledger's prevailing state --- the door never assumes a value is in the
current frame merely because it arrived now. \texttt{ingest} also refuses a
stamp naming an unregistered unit, so the condition that every stamped unit be
registered (Invariant~\ref{inv:basis}) is discharged at the door, never at
consumption. Sourcing remains the
market-data satellite's concern; the stamp and the coherence check are
ledger-owned.

The stamp assignment rests on a named trust assumption:

\begin{description}[style=nextline]
\item[TA-BASIS.] \emph{Owner:} market-data operations. \emph{Content:} the per-(unit,
source) basis stamp assigned at ingest is true of the source's dissemination.
\emph{Violation consequence:} mis-based valuation. \emph{Detection:} W3 partition
quarantine (\S\ref{sec:basis-failure}) --- a large innovation coinciding with no
committed boundary quarantines the series --- and daily reconciliation of
vendor-published adjustment factors against applied declarations.
\end{description}

\subsection{The market-data contract}
\label{sec:basis-contract}

The boundary chain has one author: the ledger. Every component outside the
ledger reads the chain and never writes it. The contract below keeps that so.

The ledger publishes one read-only value, the basis projection: a versioned
map from each registered unit to its prevailing basis id, with the unit's
committed boundary chain behind it. The market-data layer holds that value and
presents every raw number at the single ingest door
(\S\ref{sec:basis-ingest}). There the ledger attaches the stamp --- the record
of which frame the value inhabits --- from the pinned projection and the
source's declared convention (Figure~\ref{fig:md-layers}). Two flows and
nothing else: the projection flows out, observations flow in, and no basis
coordinate is authored anywhere but inside the ledger.

The market-data layer is deliberately passive --- no per-datum query, no callback,
no re-implemented corporate-action arithmetic, which the arithmetic-free projection
makes impossible even by mistake; the pinned view is an immutable value, so a
boundary committing mid-run changes nothing the layer holds, and a source lagging
the chain --- the pre-split publisher above --- is not an error to be chased down
but data in the convention on file, stamped as such.

\begin{principle}[Stamping authority]
\label{prin:stamping-authority}
The ledger's committed boundary chain is the single source of basis truth: the
ingestion layer stamps by consulting a derived, read-only projection of that chain,
never by re-implementing corporate-action logic, and no component outside the ledger
authors, defaults, or adjusts a basis coordinate.
\end{principle}

\begin{figure}[tb]
\centering
\begin{tikzpicture}[
    >={Stealth[length=2.4mm]},
    layer/.style={draw, rounded corners, align=center, font=\small,
                  minimum width=92mm, minimum height=12mm, inner sep=3pt},
    core/.style={draw, rounded corners, align=center, font=\small, fill=black!5,
                 minimum width=92mm, minimum height=12mm, inner sep=3pt},
    lbl/.style={font=\scriptsize, align=center},
    flow/.style={->, thick},
  ]
  % --- four layers, external at the top, authoritative core at the base ---
  \node[layer] (feeds)  at (0,0)     {\textbf{Feeds}\\ \scriptsize real-time / batch / historical sources, signed};
  \node[layer] (enrich) at (0,-1.7)  {\textbf{Enrichment / stamping layer}\\ \scriptsize transcription only --- model-free, arithmetic-free};
  \node[layer] (door)   at (0,-3.7)  {\textbf{Ledger ingest door} \texttt{ingest}\\ \scriptsize the sole door; the ledger attaches the stamp};
  \node[core]  (state)  at (0,-5.4)  {\textbf{Internal state}\\ \scriptsize \texttt{UnitStatus} (\texttt{usBasis}) + committed chain --- single source of basis truth};

  % --- the ledger boundary: the door and the core lie inside it ---
  \draw[dashed] (-5.2,-2.7) -- (5.2,-2.7);
  \node[lbl, fill=white] at (0,-2.7) {ledger boundary};

  % --- exactly two arrows cross the boundary ---
  \draw[flow] (-2.4,-2.3) -- node[lbl, left]{\texttt{RawDatum}\\ then \texttt{StampedObs}} (-2.4,-3.1);
  \draw[flow] (2.4,-3.1) -- node[lbl, right]{\texttt{BasisView}\\ projection (read-only)} (2.4,-2.3);
\end{tikzpicture}
\caption{The four layers as an architecture. External feeds and the enrichment/stamping
layer sit outside the ledger; the ingest door and the internal state sit inside. The single
source of basis truth is the committed boundary chain, materialised in \texttt{UnitStatus} at
the base (\S\ref{sec:basis-ingest}, Principle~\ref{prin:stamping-authority}); every layer
above reads a derived, read-only projection of it and none authors, defaults, or adjusts a
basis coordinate. Only two arrows cross the ledger boundary: raw observations descend through
\texttt{ingest} and return as stamped, while the \texttt{BasisView} projection ascends.}
\label{fig:md-layers}
\end{figure}

\paragraph{The basis projection.} One interface serves every ingestion layer ---
in-process, gateway, batch loader --- and diagnostics. \texttt{basisView} projects a
\texttt{BasisView} value from the ledger. \texttt{viewAsOf} constructs the view as of
any log position (the P8 machinery). \texttt{betaAt} returns a unit's basis id; it
returns \texttt{Nothing} exactly when the unit is unregistered, so the door's refusal
is a pattern match that cannot be forgotten. \texttt{chainAt} returns the unit's
boundary chain, ascending in the effective order of Definition~\ref{def:basis-id}.
\texttt{onChain} answers whether a claimed basis id is a committed point of the
unit's chain. \texttt{kindAt} is the registry's \texttt{betaAt}: it returns a
registered datum kind's entry from the projection of \S\ref{sec:basis-kinds}, pinned
by the same view version. Each view carries its version: the as-of log position.
Nothing else is exposed.

\begin{lstlisting}[language=Haskell]
data BasisView    -- ABSTRACT: a READ-ONLY projection of the log, same rank as
                  -- UnitStatus -- derivable, discardable, never authoritative.
                  -- Immutable once constructed; monotone under append (P24
                  -- corollary: the tip never regresses); arithmetic-free -- no
                  -- operator kinds, no parameters, no arrows are exposed.

basisView   :: Ledger -> BasisView                       -- the projection
viewAsOf    :: [Transaction] -> Timestamp -> BasisView   -- the P8 clone machinery
viewVersion :: BasisView -> ViewVersion                  -- the pin a stamp records

-- | Nothing = the unit is not registered (the door refuses, O5); Just b = the
--   prevailing basis (replayed usBasis; by P24, the committed chain's tip).
betaAt :: BasisView -> UnitId -> Maybe BasisId

-- | The unit's committed chain, ascending in the effective order (t_eff, prec,
--   bid) -- the read a re-stamp repair or a backward transport pins against.
chainAt :: BasisView -> UnitId -> [(Timestamp, Integer, BoundaryId)]

-- | The admission predicate: is b a committed point of u's chain? The origin
--   of a REGISTERED unit is the chain's base point; a boundary id is on-chain
--   iff the chain carries it. A claim naming anything else is refused
--   (OffChainClaim) -- in particular a pro-forma claim on a not-yet-committed
--   boundary, which is re-presented after the boundary commits at effectiveness.
onChain :: BasisView -> UnitId -> BasisId -> Bool

-- | The registry's betaAt (sec:basis-kinds): Nothing = the kind is not in the
--   pinned vocabulary, and the door refuses it as its FIRST clause.
kindAt :: BasisView -> KindId -> Maybe DatumKind

data Convention
  = TracksChain          -- caught up: the value is in the chain position at t_obs
  | LagsBy Integer       -- still publishing k boundaries behind that position
  | PreAdjusts           -- adjusts at commit: the value is at the pinned view's tip
  | DeclaredAt BasisId   -- file declared adjusted-through a named basis (O6);
                         --   the door checks the claim against the chain
  | NoClaim              -- nothing on file: refused, quarantined (O5)
  deriving (Eq, Show)

data RawDatum = RawDatum
  { rdUnitRef :: String                     -- provider's unit reference (O1)
  , rdValue   :: Integer                    -- exact disseminated value (O1)
  , rdConv    :: (ConvVersion, Convention)  -- convention consulted (O2/O4)
  , rdKind    :: KindId                     -- the claimed datum kind: a NAME resolved
                                            --   against the registry at the door, never
                                            --   a declaration -- the declaration lives
                                            --   on the registered entry (sec:basis-kinds)
  } deriving (Eq, Show)

data StampedObs   -- ABSTRACT: constructor unexported; ingest is the only door.
                  -- Carries value, t_obs, source, stamp :: Map UnitId BasisId
                  -- (singleton for a plain observation), and the envelope: the
                  -- projection version and convention version consulted (O2),
                  -- so every stamp replays as a pure function of recorded data.

-- The Left branch IS the quarantine: refused with payload retained, never
-- given a guessed basis, never defaulted from the ledger's prevailing state.
data IngestError
  = UnregisteredUnit  UnitId RawDatum          -- unresolvable reference: refused outright
  | UndeterminedBasis UnitId RawDatum          -- no convention on file: quarantined
  | OffChainClaim     UnitId BasisId RawDatum  -- claim names no committed chain point
  | UnregisteredKind  KindId RawDatum          -- claimed kind not in the pinned registry
                                               --   (sec:basis-kinds): the error names the
                                               --   missing REGISTRATION, and the guard is
                                               --   the door's first clause
  deriving (Eq, Show)

retained :: IngestError -> RawDatum

-- | Stamp one raw observation against one pinned view. Total; deterministic in
--   its visible arguments (pDet). The door refuses (Left, payload retained)
--   rather than guess.
ingestAt :: BasisView -> SourceId -> Timestamp -> RawDatum
         -> Either IngestError StampedObs

-- | The published door: the sole path from a raw vendor number to a
--   consumable observation -- the parse-boundary idiom of `move`. The
--   market-data layer links against `ingestAt` only: the Ledger value never
--   crosses the boundary; the projection does.
ingest :: Ledger -> SourceId -> Timestamp -> RawDatum
       -> Either IngestError StampedObs
ingest l = ingestAt (basisView l)
\end{lstlisting}

\noindent
The exhibit is drawn from the reference (\texttt{reference/Ledger.hs}, Part~M).
The signatures, \texttt{Convention}, \texttt{RawDatum}, \texttt{IngestError}, and
the \texttt{ingest} factoring are verbatim, under the code-level claim of
\S\ref{sec:states-reference}: tokens exact, comments and layout re-set for the
page. \texttt{BasisView} and \texttt{StampedObs} appear in interface form ---
declaration plus stated contract, constructors withheld, the idiom of
\S\ref{sec:basis-ingest} --- because the reference exports neither constructor.
For an abstract type the interface is the specification, and the record body is
implementation.

\texttt{BasisView} is a projection of the log like every other: it has the rank
of \texttt{unitStatus}, under the cache discipline of
\S\ref{sec:states-unitstatus} --- derivable, discardable, never authoritative.
Three of its facts carry the contract. First, it is an immutable value, so a
stamp computed from it is a pure function of its arguments. Second, it is
monotone under append, a corollary of P24: the tip never regresses, so an
on-chain stamp stays on-chain under every later view. Third, it is
arithmetic-free: it exposes no operator kinds, no parameters, no conversions.

The published door keeps the signature of \S\ref{sec:basis-ingest}, factoring as
$\texttt{ingest}\ l=\texttt{ingestAt}\ (\texttt{basisView}\ l)$. The market-data
layer links against \texttt{ingestAt} alone, so the \texttt{Ledger} value never
crosses the boundary. Three facts fix the door's shape: \texttt{ingestAt} is a
pure, total function of a value argument; \texttt{BasisView} has no mutators;
and \texttt{StampedObs} exports no constructor. The \texttt{Left} branch is the
quarantine: \texttt{IngestError} is \texttt{UnregisteredKind},
\texttt{UnregisteredUnit}, \texttt{UndeterminedBasis}, or \texttt{OffChainClaim} ---
the door's refusal order --- each retaining the
refused payload, and the quarantine store is the accumulation of the
\texttt{Left}s, inspected under the W3 workflow (\S\ref{sec:basis-failure}).
\emph{Consumable} is therefore a type fact: an observation is consumable exactly
when it is in the \texttt{Right} image of \texttt{ingest}. No unstamped
observation is consumable, by that fact together with door soundness (P25,
below) and \hyperref[cond:C13]{C13} at the seam. The door is idempotent by
content address: a duplicate raw datum returns the existing \texttt{StampedObs}.
This is the same idempotence P6 gives moves, here at the observation plane.

\paragraph{Direction of flow.} Exactly two data flows cross the market-data
boundary (Figure~\ref{fig:md-flow}). The basis projection flows out: read-only,
versioned, a view of the log. Raw observations flow in through \texttt{ingest},
the sole door, where the ledger itself attaches the stamp. The market-data layer
never writes ledger state, and the ledger never accepts a basis it did not stamp
or verify against its own committed chain. The traffic is pull-based: the layer
pulls the projection when it opens a run and pushes observations through the
door, and there is no per-datum traffic in either direction.

\begin{figure}[tb]
\centering
\begin{tikzpicture}[
    >={Stealth[length=2.4mm]},
    box/.style={draw, rounded corners, align=center, font=\small,
                minimum width=26mm, minimum height=13mm, inner sep=3pt},
    door/.style={draw, thick, rounded corners, align=center, font=\small,
                 minimum width=26mm, minimum height=13mm, inner sep=3pt},
    flow/.style={->, thick},
    lbl/.style={font=\scriptsize, align=center},
  ]
  % --- outside the ledger: feeds and the stamping layer ---
  \node[box]  (feeds)  at (0,0)    {Feeds\\ \scriptsize real-time / batch /\\ historical, signed};
  \node[box]  (enrich) at (4.3,0)  {Enrichment /\\ stamping layer\\{\scriptsize (transcription)}};
  % --- inside the ledger ---
  \node[door] (door)   at (8.8,0)  {Ledger ingest\\ door \texttt{ingest}};
  \node[box]  (state)  at (13.3,0) {Internal state\\ log + \texttt{UnitStatus}\\{\scriptsize \texttt{usBasis} + chain}};

  % --- the ledger trust boundary: the door and the state lie inside it ---
  \draw[dashed, rounded corners] (6.9,-1.15) rectangle (14.9,1.15);
  \node[lbl, anchor=south west] at (6.95,1.2) {ledger boundary};

  % --- raw observations flow in; the door attaches the stamp ---
  \draw[flow] (feeds)  -- node[lbl, above]{\texttt{RawDatum}} (enrich);
  \draw[flow] (enrich) -- node[lbl, above]{\texttt{RawDatum}} (door);
  \draw[flow] (door)   -- node[lbl, above]{\texttt{StampedObs}} (state);

  % --- the basis projection flows outward, read-only, versioned ---
  \draw[flow] (door.north) .. controls +(0,1.5) and +(0,1.5) ..
    node[lbl, above]{\texttt{BasisView} projection\\{\scriptsize read-only, versioned}} (enrich.north);
\end{tikzpicture}
\caption{The market-data contract as a data flow. Real-time and batch feeds hand raw
vendor numbers to the external enrichment/stamping layer, whose whole competence is
transcription---of values, of times, of conventions---never adjustment. Exactly two arrows
cross the ledger boundary (\S\ref{sec:basis-contract}, Principle~\ref{prin:stamping-authority}):
the read-only, versioned \texttt{BasisView} projection flows out, and raw observations flow
in through \texttt{ingest}, the sole door, where the ledger itself attaches the basis stamp
and commits the accepted observation as a logged event; the stamp is checked against
\texttt{UnitStatus} (\texttt{usBasis},
Definition~\ref{def:usbasis}) over the committed boundary chain. The market-data layer never
writes ledger state; the ledger never accepts a basis it did not stamp or verify against its
own committed chain.}
\label{fig:md-flow}
\end{figure}

\paragraph{The obligation set.}
\begin{description}[style=nextline]
\item[O1 (raw observation).] A raw observation carries exactly: the value, the
observation time $t_{\mathrm{obs}}$, a signed source id, and a unit reference
resolvable to a registered \texttt{UnitId}. Nothing more is demanded of a provider.
$t_{\mathrm{obs}}$ is attested data; it is never compared against a wall clock to
decide semantics.
\item[O2 (the stamp).] The ledger attaches the stamp at ingestion, from the pinned
basis projection and the source basis convention (O4); the attestation envelope
records the projection version and the convention version consulted. Every stamp
therefore replays as a pure function of recorded data: same projection version,
convention version, source, $t_{\mathrm{obs}}$, and raw datum give a bit-identical
\texttt{StampedObs}.
\item[O3 (no corporate-action awareness).] A provider is not required to track
corporate actions: a pre-adjusting source and a lagging source are both
representable, because the stamp states which basis the value is in --- the quote-ex
flag of \S\ref{app:pricing-coordination} generalised from one bit to the coordinate.
\item[O4 (source basis convention).] Each (unit, source) pair has a \emph{source
basis convention} on file: how the source's dissemination relates to the chain ---
tracks it, lags by a declared offset, pre-adjusts, declares a fixed basis, or makes
no claim. Conventions enter as logged observation events under the canonical-handler
idiom of Definition~\ref{def:boundary-event}: versioned, as-of-queryable, no new
state home. Owner: market-data operations, under TA-BASIS
(\S\ref{sec:basis-ingest}).
\item[O5 (fail-closed).] A datum whose basis the convention and the pinned view
cannot determine is refused with payload retained --- never given a guessed basis,
never defaulted from the ledger's prevailing state; an unresolvable or unregistered
unit reference is refused outright (\S\ref{sec:basis-ingest}).
\item[O6 (back-adjusted files).] A historical series pulled as a back-adjusted file
is stamped from the file's declared adjusted-through convention, never from the
basis prevailing at each row's observation time.
\item[O7 (signature).] Feeds are signed; an unsigned feed enters only through a
gateway that counter-signs under a named key, carrying its own scoped trust
assumption.
\item[O8 (corrections).] A vendor value correction is a new observation, superseding
by (source, unit, $t_{\mathrm{obs}}$); a duplicate content address is idempotent at
the door; basis re-coordination is exclusively the ledger's re-stamp event with
lineage (\S\ref{sec:basis-time-travel}).
\end{description}

\paragraph{Primitive and derived.} A datum is \emph{derived} exactly when it is
computed inside a snapshot scope from ledger-stamped inputs; it carries the joint
stamp of its inputs and is never transported (Principle~\ref{prin:rederivation}).
A value computed outside the ledger from data the ledger never stamped --- a
provider-published index level, a vendor-supplied curve --- is a \emph{primitive}
observation of the publisher's unit, with a singleton stamp; the publisher's
rebalance and divisor events enter as W4 basis notices on that unit. The dividing
line is that authority follows computation: the log has no authority over
arithmetic it did not perform, and a joint stamp on a provider-computed composite
would attest to nothing, because the ledger never saw the inputs.

The per-kind run of this contract --- for each registered datum kind, what the stamp
covers and what fail-closed means --- is the table of \S\ref{sec:basis-kinds}, whose
registry fixes each kind's component dimensions.

\paragraph{One contract, three modes.} A real-time feed, an end-of-day batch, and
a historical pull are three modes of arrival, and all three face the same door
under the same contract: mode is not among the stamp's arguments. One ingestion
run pins one projection version. Within the run, each stamp is computed by the
source basis convention from the datum and its $t_{\mathrm{obs}}$, against the
pinned chain. $t_{\mathrm{obs}}$ is an input the convention may use; it is never
the coordinate itself. Where the convention names the basis directly --- a
pre-adjusting source, or a file declared adjusted-through (O6) --- that
declaration overrides any positional reading off the chain at
$t_{\mathrm{obs}}$, and the stamp is never defaulted from the basis prevailing
at arrival (\S\ref{sec:basis-ingest}). A boundary committing mid-run cannot tear
the run, because the run's view is one immutable value. Batch and real-time
therefore agree by construction: the stamp is a pure function of projection
version, convention version, source, $t_{\mathrm{obs}}$, and raw datum. Multiple
vendors are the same door, one per \texttt{SourceId}. Two vendors disagreeing on
value within one basis partition is a data-quality question, outside this
contract; two vendors disagreeing on basis is two different stamps, resolved by
W3's partitions (\S\ref{sec:basis-failure}), across which no aggregate can be
computed.

\paragraph{Two consumers across one boundary.} The 2-for-1 split is effective at
the market open. The stored pre-split print, 100, is stamped $b_3$ (the unit's basis before the
split; the subscript is just its position in this unit's boundary chain, which began
before the split). Consumer A
calls \texttt{withSnapshot} --- the operation that opens a read scope fixed to
one basis assignment, defined fully at the typed seam
(\S\ref{sec:basis-types}) --- at 09:29. In A's scope $\beta(u)=b_3$: the print is
consumed as stored, and the mark is $1{,}000\times100=\EUR{100{,}000}$. Consumer
B opens its scope at 09:31. In B's scope $\beta(u)=b_4$: the same stored
observation is transported along the unique declared conversion to 50, and the
mark is $2{,}000\times50=\EUR{100{,}000}$. One stored series, read from opposite
sides of the boundary, gives two correct marks, and the totals agree by
Theorem~\ref{thm:basis-value-invariance} with per-share cash $c=0$. The phantom
$2{,}000\times100$ can be formed in neither scope: A's snapshot holds $1{,}000$
shares, and in B's scope the raw 100 has no way in --- \texttt{snapPrice} and
\texttt{transport} yield 50 or refuse. Neither consumer knows the split's terms;
each knows only its own $\beta$ and the door.

No corporate-action logic exists outside the ledger, and none can be built
there: the projection is arithmetic-free, so an external component does not hold
the information adjustment would need. Adjustment conversions are declared once,
in the log, and applied by the one evaluator. The ingestion layer's whole
competence is transcription --- of values, of times, of conventions --- never
adjustment.

\paragraph{Boundary faults.} The consequence classes are the four regimes of
\S\ref{sec:basis-failure}. The table below lists the fault scenarios that arise
at this boundary and where each one discharges.

\renewcommand{\arraystretch}{1.15}
\begin{center}
\begin{longtable}{@{}p{4.0cm} p{9.4cm}@{}}
\toprule
\textbf{Fault} & \textbf{Discharge} \\
\midrule
\endhead
F1 Late boundary notice & Tip weld (Principle~\ref{prin:tip-weld}); re-stamp events
with lineage (\S\ref{sec:basis-time-travel}); W3 detects and W1 blocks in the
window. \\
F2 Duplicated notice & Content address; \texttt{SetBasis} idempotent (P6). \\
F3 Reordered notices & Every admissible booking permutation yields the identical
chain projection --- the tip weld again; property form P-PERM-N. \\
F4 Omitted notice & W3 innovation-without-boundary quarantine --- detection, not
guarantee; the window is fail-closed. \\
F5 Duplicated observation & The door is idempotent by content address and returns
the existing \texttt{StampedObs}. \\
F6 Reordered or late observations & Valuation is invariant under ingestion-order
permutation at fixed pins; property form P-PERM-O. \\
F7 Omitted observation & \texttt{snapPrice} returns \texttt{Left};
\texttt{Unpriced} enumerated (W2). \\
F8 Clock skew & $t_{\mathrm{obs}}$ is data (O1); the only clock read is the W1
staleness check, whose threshold lives in \emph{ProductTerms} and whose clock is an
explicit invocation argument, never ambient. \\
F9 Crash mid-ingestion & A \texttt{StampedObs} exists only as a logged event: a
crash before append is nothing-happened, resubmission dedups (F5), and recovery
reaches a committed prefix; property form P-CRASH. \\
F10 Multi-vendor basis disagreement & W3 partitions by stamped basis; no
cross-partition aggregate is representable --- the property form of the W3 rule. \\
\bottomrule
\end{longtable}
\end{center}

\paragraph{Property forms.} The boundary contributes nine property forms, each
stated once.
\begin{description}[style=nextline]
\item[P25, door soundness.] On \texttt{Right}, every stamped unit is registered and
every stamped basis id is on the unit's chain; on \texttt{Left}, nothing is admitted
and the payload is retained.
\item[P-DET, stamp determinism.] The stamp is a pure function of recorded data ---
the property form of O2.
\item[P-MODE, mode equivalence.] Mode is not an argument of the stamp: real-time,
batch, and historical arrival produce the identical stamp from identical inputs.
\item[P-PERM-N, P-PERM-O.] Permutation invariance: booking the same notices in any
admissible order yields the identical chain, and ingesting the same observations in
any order at fixed pins yields the identical valuations.
\item[P-CRASH.] Recovery reaches a committed prefix.
\item[P-REPRO.] Same log, same stamped observations, same pinned chain version ---
bit-identical valuations.
\item[P-CLONE-STAMP.] The as-known replay reproduces the honest mistake; the corrected
replay applies the re-stamps; both replays are deterministic (\S\ref{sec:basis-time-travel}).
\item[P-LAG, lag arithmetic.] A \texttt{LagsBy}~$k$ stamp sits exactly $k$
effective-order steps behind the caught-up position at the same $t_{\mathrm{obs}}$,
floored at the origin. This catches an off-by-one that door soundness cannot see,
because any on-chain id is sound.
\end{description}
The executable forms follow, verbatim from the reference (\texttt{reference/Ledger.hs},
Part~M), in the catalogue shape of Appendix~\ref{app:properties}. A claim about a single
invocation (P25) is written as a property, following P1 and P2; a claim that quantifies
over a whole log or feed is written as a boolean oracle --- a function that returns
true exactly when the property holds --- following P8 and P24.

\begin{lstlisting}[language=Haskell]
-- P25 Ingest-door soundness. Run o = ingest l s t rd. On Right: every stamped
-- unit is registered and every stamped id is a committed point of that unit's
-- chain (origin included) -- nothing enters the consumable plane claiming a
-- basis the ledger has not committed. On Left: nothing is admitted and the
-- payload is retained verbatim -- the Left IS the quarantine. Wrongful-
-- rejection is the vacuous half, as in P1/P2. The no-side-door half is the
-- type fact stated at StampedObs; the consumption half is C13 at withSnapshot.
p25 :: Property (SourceId, Timestamp, RawDatum) (Either IngestError StampedObs)
p25 = Property (\_ _ -> True) $ \l (_, _, rd) o -> case o of
  Left e   -> retained e == rd
  Right so -> all (\(u, b) -> registered l u && onChain (basisView l) u b)
                  (Map.toList (obsStamp so))

-- P-DET Stamp determinism, restated as the purity witness (the pureL
-- precedent): the stamp is a total pure function of (projection version,
-- convention version, source, t_obs, raw) and of NOTHING else -- every input
-- is a value argument; no clock, no ambient log read.
pDet :: BasisView -> SourceId -> Timestamp -> RawDatum -> Bool
pDet v s t rd = ingestAt v s t rd == ingestAt v s t rd

-- P-MODE Mode equivalence -- the executable shadow of a one-sentence theorem:
-- the stamp is a pure function of the pinned view and the datum, so arrival
-- MODE (tick-by-tick, end-of-day file, resumed run) cannot enter the result.
-- One datum through the door equals the same datum inside any batch at the
-- same pin; a boundary committing mid-run cannot tear the run, because the
-- run never re-reads the ledger.
pMode :: BasisView -> [(SourceId, Timestamp, RawDatum)] -> Bool
pMode v feed =
  fst (ingestRun v feed)
    == Set.fromList [ so | (s, t, rd) <- feed, Right so <- [ingestAt v s t rd] ]

-- P-PERM-N Notice-order irrelevance, strong form: the effective order is data
-- in the notices ((t_eff, prec, bid)), so booking order cannot enter the chain
-- projection or the tip. A naive permutation of a committed log is mostly
-- REFUSED by the tip weld -- each SetBasis re-assertion is derived for its own
-- booking position -- and a refused replay asserts nothing (P2's half), so the
-- naive trial is vacuous almost always. The oracle therefore re-derives what
-- is derived and permutes only what is data: the boundary NOTICES are permuted
-- among their own positions (registrations stay put, permuteBoundariesBy) and
-- each SetBasis is re-asserted to the post-insertion tip of the NEW order
-- (rebook), exactly as genChain mints it. Every permutation of a generated
-- chain then replays Right, and the trial is non-vacuous by construction; the
-- vacuous branch remains only for adversarial inputs, where refusal is P2's
-- half, as before.
pPermN :: [Transaction] -> [Int] -> Bool
pPermN txs perm =
  case ( replay txs emptyLedger
       , replay (rebook (permuteBoundariesBy perm txs)) emptyLedger ) of
    (Right a, Right b) ->
      all (\u -> chainAt (basisView a) u == chainAt (basisView b) u
              && betaAt  (basisView a) u == betaAt  (basisView b) u)
          (map txUnit txs)
    _ -> True

-- P-PERM-O Observation-order irrelevance: at a fixed pin, the committed set of
-- a run is invariant under any permutation of arrivals -- each datum's stamp
-- is a function of the pin and the datum, never of its neighbours. (The
-- quarantine LIST keeps arrival order and is not claimed; the committed SET is.)
pPermO :: BasisView -> [(SourceId, Timestamp, RawDatum)] -> [Int] -> Bool
pPermO v feed perm =
  fst (ingestRun v feed) == fst (ingestRun v (permuteBy perm feed))

-- P-CRASH Crash recovery: a run interrupted after k data leaves exactly its
-- committed prefix -- a StampedObs exists only as an appended event, so
-- crash-before-append is nothing-happened -- and resubmitting the WHOLE feed
-- against the SAME pin dedups by content address and lands on the
-- uninterrupted run's committed set. Every crash point recovers to a committed
-- prefix, and then to the run itself.
pCrash :: BasisView -> [(SourceId, Timestamp, RawDatum)] -> Int -> Bool
pCrash v feed k =
  fst (ingestResume v (fst (ingestRun v (take k feed))) feed)
    == fst (ingestRun v feed)

-- P-REPRO Reproducibility (as-of stability): the view as-of booking position c
-- is unchanged by anything booked after c -- viewAsOf is prefix-stable -- so a
-- stamp replayed today against the recorded pin is bit-identical to the stamp
-- minted then. Same log, same observations, same pinned version: the same
-- StampedObs, and (with P8 and C13) the same valuations.
pRepro :: [Transaction] -> Timestamp -> Int -> (SourceId, Timestamp, RawDatum) -> Bool
pRepro txs (Timestamp c) ext (s, t, rd) =
  let m = fromInteger (max 0 c) + max 0 ext          -- any later extension
  in ingestAt (viewAsOf (take m txs) (Timestamp c)) s t rd
       == ingestAt (viewAsOf txs (Timestamp c)) s t rd

-- P-CLONE-STAMP clone_at stamp reconstruction: for a stamp minted at ANY
-- booking cut, the log alone reconstructs the pin -- some cut's rebuilt view
-- carries the recorded projection version (a content address, so a
-- retro-effective boundary booked later renumbers nothing) -- and EVERY cut
-- carrying that version re-mints the stamp bit-identically. The as-known-at-t
-- replay thus reproduces the honest mistake rather than repairing it; the
-- corrected reading is the same door at the tip pin (deterministic by pDet);
-- re-coordination beyond that is exclusively the ledger's re-stamp event with
-- lineage (O8), outside this oracle.
pCloneStamp :: [Transaction] -> (SourceId, Timestamp, RawDatum) -> Bool
pCloneStamp txs (s, t, rd) = all cutOK cuts
  where
    cuts = [0 .. toInteger (length txs)]
    viewAt k = viewAsOf txs (Timestamp k)
    cutOK k = case ingestAt (viewAt k) s t rd of
      Left _   -> True                     -- a refusal asserts nothing (P2's half)
      Right so ->
        let hits = [ j | j <- cuts, bvVersion (viewAt j) == obsView so ]
        in  not (null hits)
            && all (\j -> ingestAt (viewAt j) s t rd == Right so) hits

-- The convention's positional readings, internal to the door: the chain
-- position at t in EFFECTIVE order (the maximal committed boundary with
-- t_eff <= t, or the origin), and the position k boundaries behind it,
-- floored at the origin. Total; consumed by TracksChain / LagsBy; never a
-- stamp default -- a datum with no convention is refused, never positioned.
chainPosAt :: BasisView -> UnitId -> Timestamp -> BasisId
lagPosAt   :: BasisView -> UnitId -> Timestamp -> Integer -> BasisId

-- P-LAG Lag-convention arithmetic: a LagsBy k stamp sits exactly k
-- effective-order steps behind the caught-up position at the same t_obs,
-- floored at the origin -- and LagsBy 0 IS TracksChain. P25 cannot see an
-- off-by-one here, because any on-chain id is sound; the arithmetic therefore
-- carries its own oracle, computed independently by positional indexing where
-- lagPosAt drops from the reversed chain. Drivers: views from genChain via
-- viewAsOf, t and k drawn by nextInt, k shrunk by shrinkInteger.
pLag :: BasisView -> UnitId -> Timestamp -> Integer -> Bool
pLag v u t k =
  let path = [ bid | (te, _, bid) <- chainAt v u, te <= t ]   -- ascending
      k'   = fromInteger (max 0 k)
      want | length path <= k' = Origin u
           | otherwise         = Boundary (path !! (length path - 1 - k'))
  in lagPosAt v u t k == want  &&  lagPosAt v u t 0 == chainPosAt v u t

-- F10 Basis-partition separation: observations aggregate only within one stamp
-- cell. The oracle witnesses the partition's soundness; that no AGGREGATE can
-- consume two cells is the consumption seam's type fact (C13), not re-proved
-- here.
pPartition :: [StampedObs] -> Bool
pPartition sos =
  all (\(st, cell) -> all ((== st) . obsStamp) cell)
      (Map.toList (partitionByStamp sos))
\end{lstlisting}

\noindent
The generators and shrinkers that drive the catalogue are deterministic and pure,
with no external dependency: a \texttt{Seed} drives a splittable linear-congruential
generator (splittable so independent draws share no state),
every generator is a total function of its \texttt{Seed}, and every shrinker enumerates
purely toward the structurally smaller candidate. A full body is shown only where the
way a value is built is itself part of the argument: a \texttt{StampedObs} can be
generated only by calling the door, and a chain only by building a committed log. Where
the generator is just plain seed-driven drawing, only its signature and contract are
shown.

\begin{lstlisting}[language=Haskell]
newtype Seed = Seed Integer deriving (Eq, Show)

nextInt   :: Integer -> Seed -> (Integer, Seed)   -- uniform-ish in [0, n); total
splitSeed :: Seed -> (Seed, Seed)

-- | genVendorBehaviour: the source basis convention alphabet -- caught-up,
--   lagging, pre-adjusting, declared-through (on-chain ids AND an off-chain id,
--   so the OffChainClaim refusal is exercised), and nothing-on-file. Shrinks
--   toward TracksChain, the best-behaved vendor. The forged off-chain id is
--   harmless by construction: the door refuses it, and the tip weld refuses it
--   as a SetBasis -- it can never name a committed state.
genVendorBehaviour :: BasisView -> UnitId -> Seed -> Convention

shrinkVendorBehaviour :: Convention -> [Convention]
shrinkVendorBehaviour TracksChain = []
shrinkVendorBehaviour (LagsBy k)  = TracksChain : [ LagsBy k' | k' <- shrinkInteger k ]
shrinkVendorBehaviour _           = [TracksChain]

-- | genRawDatum: draws the unit reference from the view's register, PLUS one
--   unregistered reference so the UnregisteredUnit refusal is exercised; an
--   exact value; and a generated vendor behaviour with its convention version.
genRawDatum   :: BasisView -> Seed -> RawDatum
shrinkRawDatum :: RawDatum -> [RawDatum]
shrinkRawDatum (RawDatum ref val (cv, conv)) =
     [ RawDatum ref v' (cv, conv) | v' <- shrinkInteger val ]
  ++ [ RawDatum ref val (cv, c')  | c' <- shrinkVendorBehaviour conv ]

-- | genStampedObs: ONLY by calling the door. A generator that built a
--   StampedObs by record syntax would mint values outside ingest's image --
--   precisely the states the abstract type excludes -- so generation is
--   ingestAt over generated inputs, Nothing when the door refuses, and
--   shrinking shrinks the RAW DATUM and re-ingests (shrinkStampedObs): a
--   stamp is never shrunk in place.
genStampedObs :: BasisView -> Seed -> Maybe StampedObs
genStampedObs v sd0 =
  let (s1, s2) = splitSeed sd0
      (tv, _ ) = nextInt 16 s1
  in either (const Nothing) Just
       (ingestAt v (SourceId "gen") (Timestamp tv) (genRawDatum v s2))

shrinkStampedObs :: BasisView -> SourceId -> Timestamp -> RawDatum -> [StampedObs]
shrinkStampedObs v s t rd =
  [ so | rd' <- shrinkRawDatum rd, Right so <- [ingestAt v s t rd'] ]

-- | genFeed: a multi-vendor feed against one view -- several sources, mixed
--   vendor behaviours, observation times straddling the chain's boundaries.
--   genPerm and genCrashPoint drive the pPermO / pCrash schedules over the
--   same feed; partitionByStamp splits its committed set for pPartition.
genFeed       :: BasisView -> Int -> Seed -> [(SourceId, Timestamp, RawDatum)]
genPerm       :: Int -> Seed -> [Int]
genCrashPoint :: Int -> Seed -> Int

-- | Shrinkers, paired: a feed shrinks by prefix, then pointwise through
--   shrinkRawDatum (source and t_obs name the trial; the datum carries the
--   structure); a permutation shrinks one step, to the identity; a crash
--   point shrinks toward 0, then halfway.
shrinkFeed       :: [(SourceId, Timestamp, RawDatum)]
                 -> [[(SourceId, Timestamp, RawDatum)]]
shrinkPerm       :: [Int] -> [[Int]]
shrinkCrashPoint :: Int -> [Int]

-- | The P-PERM-N drivers: permuteBoundariesBy permutes the boundary NOTICES
--   among their own positions (registrations stay put); rebook re-asserts
--   each derived SetBasis to the post-insertion tip of the new order, per
--   unit, exactly as genChain mints it -- so every permutation of a
--   generated chain replays Right, and no trial is vacuous.
permuteBoundariesBy :: [Int] -> [Transaction] -> [Transaction]
rebook              :: [Transaction] -> [Transaction]

-- In-module chain minting (the P24 precedent). BoundaryId is a content address
-- assigned at the attested notice gateway; the oracle harness mints
-- deterministic addresses HERE, in-module, and these minters are deliberately
-- NOT exported -- no production caller gains a minting door.

-- | genBoundaryEvent: covers retro-effective t_eff (the small window makes a
--   later booking with an earlier t_eff frequent), same-t_eff precedence
--   collisions (tie broken by the content address, data in the notice), and
--   every Declaration constructor -- parameterised (with and without a cash
--   leg), Pending, Terminal.
genBoundaryEvent    :: UnitId -> Seed -> BoundaryEvent
shrinkBoundaryEvent :: UnitId -> BoundaryEvent -> [BoundaryEvent]

-- | genChain: a committed log for one unit -- registration, then n boundary
--   transactions, each SetBasis re-asserting the post-insertion effective tip
--   so the tip weld admits it (retro-effective notices included). Holder-free
--   by construction, so the invariance weld is the vacuous zero-holder case
--   and the generator exercises the CHAIN, not entitlement arithmetic; every
--   generated log replays Right. Shrinking is prefix-closed (shrinkChain):
--   every prefix of a committed log is a committed log.
genChain    :: UnitId -> Int -> Seed -> [Transaction]
shrinkChain :: [Transaction] -> [[Transaction]]
shrinkChain txs = [ take k txs | k <- [1 .. length txs - 1] ]

-- | genBoundaryStraddle: the boundary-straddle scenario -- one unit, one
--   committed boundary, and the two booking cuts on either side of it. A
--   consumer pinned at the earlier cut reads the stored series at the old
--   point; a consumer pinned at the later cut reads it transported along the
--   declared arrow; both are correct, related by value invariance, and the
--   cross-reading is refused at the seam (C13).
genBoundaryStraddle :: Seed -> ([Transaction], Timestamp, Timestamp)
\end{lstlisting}

\noindent
Properties, generators, and shrinkers are verbatim from the reference
(\texttt{reference/Ledger.hs}, Part~M), authored and verified by inspection;
execution is pending the toolchain and is not reported as passed.

\subsection{The datum-kind registry}
\label{sec:basis-kinds}

A market datum is not one number under one name. A dividend forecast lists entries,
some a cash amount per share, some a proportion of price; a vol quote carries a
strike coordinate and a vol value. Under one boundary operator the components of one
datum can have different fates, so the declaration that governs transport is made
per component, never per name.

Each consumable datum kind is therefore \emph{registered}, once, before any datum of
that kind can enter: a kind registration event fixes the kind's identity and its
\emph{component dimension declaration} --- one declared dimension per component:
price, cash, or dimensionless. A datum is an observation, never a holding, so the
quantity dimension is excluded at the registration door. The per-dimension action
then fixes
how each component crosses a boundary, and the one declaration serves both consumers
(Principle~\ref{prin:one-decl-two-uses}): the ingest door stamps against it, and the
read seam converts by it. Spot registers a single price component. The dividend
forecast registers its cash entries as prices and its proportional entries as
dimensionless. A vol quote registers its strike coordinate as a price and its vol
value as dimensionless.

The vol declaration decides a question of identity. Absolute-strike vol and
moneyness vol are two registered kinds, not one kind with a convention flag.
An absolute strike transforms as a price under \texttt{Scale} while a
moneyness ratio stands, and two quantities that transform differently under
the same operator are different kinds: the kind id says so, and no
conditional enters the transport path. Within a kind, each component carries
its own boundary behaviour. The vol value is dimensionless and is never
transformed. An absolute-strike coordinate transforms by transcription of the
exchange's published contract adjustment, never by a model. A moneyness
coordinate crosses the boundary unchanged only where the boundary notice declares
\texttt{AId}, the declared identity conversion. Identity is declared, never
assumed: the default declaration is empty, so by default transport is
refused. The quotation convention stays on the entry as
attested provenance, checkable but never branched on.

The dividend forecast shows two fates under one operator. A wallet holds 1{,}000
shares at \EUR{100}; the forecast lists \EUR{3} cash per share and a proportional
entry of 1\% of price. The 2-for-1 split declares \texttt{Scale}~$2$. The
cash component is a price, so \EUR{3} reads as \EUR{1.50} post-split, and the cash
delivered is invariant: $1{,}000\times3 = 2{,}000\times1.50 = \EUR{3{,}000}$. The
proportional component is dimensionless, so 1\% stands, and its leg is likewise
invariant: $1\%\times100\times1{,}000 = 1\%\times50\times2{,}000 = \EUR{1{,}000}$
--- one datum, two components, two fates, and total cash unchanged on both, which is
Theorem~\ref{thm:dimension-invariance} applied per component. Now misdeclare the
proportional entry as a price. \texttt{Scale} halves it to $0.5\%$, and the leg pays
$0.5\%\times50\times2{,}000 = \EUR{500}$ --- half the true entitlement, from one
wrong word in a declaration. The mandatory witness below exists to catch this
class of error.

The registry is vocabulary, and it lives where vocabulary must: in the log.
Registration and retirement are attested events under the notice-admission
rule W4 (\S\ref{sec:basis-failure}). The registry the door consults is a
projection of those events, under the cache discipline of
\S\ref{sec:states-unitstatus}: derivable, discardable, never itself
authoritative. Each ingestion run consults the registry through the same
pinned view it stamps against --- the snapshot fixed when the run opens
(\S\ref{sec:basis-contract}) --- so every stamp and every refusal replays as
a pure function of recorded data (O2). A registration is immutable: a wrong
declaration is corrected by registering a new kind and logging a retirement
that closes the old kind to new ingestion, never by editing an entry.

An unregistered kind is unconsumable. The door refuses it as
\texttt{UnregisteredKind}, retaining the payload, as its first clause --- before
unit resolution, convention lookup, or claim checking. This refusal is distinct from
an undeclared conversion: an unregistered kind never enters the ledger at all, while
a registered kind that a boundary's declaration omits enters and is stopped only at
that crossing (the declared partiality of Definition~\ref{def:opspec}). The first is
a vocabulary failure at the door; the second is a conversion gap at a boundary.

Every registry entry carries an \emph{invariance witness}, and an entry without one
is inadmissible: registration demands the witness with the declaration. For a scalar
kind the witness is an instance of the invariance theorem itself
(Theorem~\ref{thm:basis-value-invariance}); a composite kind names an executable
oracle --- the dividend-list kind cites \texttt{dkDivSplitOK}, which replays a scale
boundary against the declaration, checks the delivered-cash identity above, and
fails the misdeclared variant. The declaration the witnesses pin is a named trust
assumption:

\begin{description}[style=nextline]
\item[TA-KIND.] \emph{Owner:} data governance, the W4 owner. \emph{Content:} each
registry entry's component dimension declaration is true of the kind it names.
\emph{Violation consequence:} dimensionally wrong transport --- the \EUR{500} leg
above. \emph{Detection:} the entry's invariance witness at registration, W3
partition divergence (\S\ref{sec:basis-failure}), and differential recomputation
(P27). \emph{Accepted residual:} a misdeclared kind whose data never cross a
boundary.
\end{description}

\noindent
TA-BASIS trusts the stamp; TA-KIND trusts the declaration of what the stamp covers.
The scope of that trust is explicit: each registered kind is trusted only to
the extent that its per-entry invariance witness holds, and a kind without
its witness cannot be registered.

The registry classifies what the ledger transports; it grants no authority over
arithmetic the ledger did not perform. A provider-computed composite --- an index
level, a vendor curve --- registers as a primitive kind of the publisher's unit,
never as a composite of its inputs, and the re-derivation canon
(Principle~\ref{prin:rederivation}) stands: derived data are re-derived from
transported inputs, never transported. Boundary action on registered kinds stays
within the closed operator menu of Definition~\ref{def:opspec}: an index publisher's
rebalance or divisor change enters as a W4 basis notice declaring
\texttt{Recompose}, and a benchmark cessation or succession is \texttt{Subst}
composed with \texttt{Shift} at the published fallback spread --- instances of
existing operators, not new cases.

The table below runs the contract across the registered kinds: for each kind, what
the stamp covers and what fail-closed means for that kind.

\renewcommand{\arraystretch}{1.15}
\begin{center}
\begin{longtable}{@{}p{2.4cm} p{5.6cm} p{5.5cm}@{}}
\toprule
\textbf{Raw datum} & \textbf{What the stamp covers} & \textbf{Fail-closed meaning} \\
\midrule
\endhead
Spot, quote & Disseminated trade or quote; singleton stamp $\{u\mapsto b\}$ per
(unit, source) convention; a lagging source is the per-source convention case
(\S\ref{sec:basis-ingest}). & No
convention on file: the series is refused into quarantine; around ex dates W3
partitions by stamped basis. \\
\addlinespace
Implied-vol observation & The vol value with its labelled coordinates --- absolute
strike, expiry, moneyness or delta --- transcribed from the source, never inferred;
singleton stamp on the quoted unit (the option series where registered, else the
underlying). No parametrisation is mandated.
& An undeclared (coordinate, operator) pair has no conversion: the point cannot cross the
boundary; an unpriced vol point is an enumerated \texttt{Unpriced}, never a
transported value. \\
\addlinespace
Index level & Provider-published level: a primitive observation of the index unit
$I$; singleton stamp $\{I\mapsto b_I\}$. The divisor snapped in-scope from stamped
constituents is the derived case and keeps its joint stamp
(\S\ref{sec:basis-worked}). & A composition change
with no committed maintenance boundary is a large innovation without a boundary: W3
quarantines the level series. \\
\addlinespace
Composition, weights, divisor files & The publisher's declared parameters for a
maintenance boundary (\texttt{Recompose}\,$\sigma$), with its $t_{\mathrm{eff}}$;
logged as a W4 notice. &
Pro-forma data claiming a basis id not yet on the committed chain are refused
(\texttt{OffChainClaim}) and re-presented after the boundary commits at
effectiveness; a stale divisor consumed across the boundary is the refused 70.83
phantom (\S\ref{sec:basis-worked}). \\
\addlinespace
Fixing, benchmark & Value, fixing time, publication time; singleton stamp on the
benchmark unit at the basis the administrator's convention names. & A republication
within the
correction window is a new observation on the value axis, basis untouched; a fixing
meeting \texttt{Terminal} or a \texttt{Pending} gap blocks under W1 --- real money
never moves on an unwitnessed basis. \\
\addlinespace
Curve, derived data & Derived in-scope from transported primitives; the joint stamp
of its inputs, inherited compositionally (Principle~\ref{prin:rederivation},
Proposition~\ref{prop:composition-law}); never transported --- the inputs transport
and the derivation re-runs at the target assignment. A vendor-supplied curve
is the primitive case of the index-level row. & One refusing input refuses the
derived point, enumerated per unit. \\
\bottomrule
\end{longtable}
\end{center}

\subsection{The market-data integrator}
\label{sec:basis-integrator}

An integrator is the component that brings external market data --- prices,
implied-volatility points, index levels --- to the ingest door. It does four
things, and nothing else.

First, it \emph{pins}. At the start of a run it takes one \texttt{BasisView}
and holds that view for the whole run. The view is immutable, so a boundary
that commits in the middle of the run cannot change what the run sees: every
datum in the run is judged against the same picture.

Second, it \emph{transcribes}. From the feed it copies the disseminated value,
the observation time, the signed source id, and the provider's unit reference
into a \texttt{RawDatum}, under the (unit, source) convention on file (O4).
Copying values, times, and conventions is the whole of its competence; it never
adjusts a number.

Third, it \emph{presents}. It passes each datum through \texttt{ingestAt}
against the pinned view. The integrator links against \texttt{ingestAt} alone,
and the \texttt{Ledger} value never crosses the boundary. The ledger attaches
the stamp; the integrator never computes a basis id.

Fourth, it \emph{keeps the refusals}. Each \texttt{Left} it receives retains
its payload, and the quarantine store --- the accumulated record of those
refusals --- is inspected under the W3 workflow
(\S\ref{sec:basis-failure}).

The integrator therefore carries no live picture of the ledger. The stamp is a
pure function of projection version, convention version, source,
$t_{\mathrm{obs}}$, and raw datum (O2, P-DET), so a source that lags a boundary
needs no intervention. After the 2-for-1 split creates the post-split basis, a
source that keeps publishing pre-split prints is on file as \texttt{LagsBy},
and the door stamps its print of 100 at the pre-split chain point. The 100 is
correct there; consumption at the post-split basis either transports it to 50
along the declared conversion or refuses it, so the stamped 100 can never meet
the 2,000-share balance unexamined. The \texttt{mQuoteEx} bit of
\S\ref{app:pricing-coordination} is exactly this stamp collapsed to one step,
generalised by the coordinate to boundary chains of any length.

Real time and batch share one code path. A tick-by-tick feed presents each
datum at the door as it arrives. An end-of-day file presents the same door,
datum by datum, within one run. A historical back-adjusted pull is stamped from
the file's declared adjusted-through convention (O6), never from the basis
prevailing at each row's observation time. Mode is not among the stamp's
arguments (P-MODE): a datum presented in real time and the same datum presented
inside a batch at the same pin receive the same stamp. The only operational
difference between the modes is when the integrator refreshes its pin ---
between runs, never within one.

\subsection{The invariant}
\label{sec:basis-invariant}

The discipline reduces to one statement. A valuation run works against a single
basis assignment: one basis per unit, fixed for the whole run. Every balance
the run reads and every datum it consumes must agree with that assignment, and
the agreement is checked over every unit the datum's stamp names, not only the
units being valued.

\begin{invariant}[Single-basis consumption]
\label{inv:basis}
Let $V$ be a valuation invocation over position scope $S$ at time $t$, and let
$\beta_t:\mathrm{UnitId}\to\mathrm{BasisId}$ be the total basis assignment given by
replayed \texttt{usBasis} over all registered units (total by
Definition~\ref{def:usbasis}: \texttt{defaultStatus}\,$u$ carries \texttt{Origin}\,$u$). Then:
(i)~every balance consumed by $V$ is read from $S$ at $\beta_t|_S$;
(ii)~every datum $d$ consumed by $V$ satisfies
$\mathrm{stamp}(d)(u)=\beta_t(u)$ for every $u\in\mathrm{dom}(\mathrm{stamp}(d))$ ---
the check ranges over the stamp's whole domain, whether or not $u\in S$;
(iii)~every value formed from several data is formed from data satisfying (ii) under one
and the same $\beta_t$.
\end{invariant}

\noindent
In the small world: a valuation whose assignment $\beta_t$ places the share at
the post-split basis may read the 2,000-share balance and a price stamped
post-split (50), but not the pre-split print of 100 --- clause~(ii) refuses the
stale stamp even though each number is individually correct, and the
\EUR{200{,}000} phantom of the opening never forms. Clause~(ii) also ranges
beyond the valuation scope: the level of index $I$ is stamped over the
constituents $A$ and $B$ it was computed from, so a valuation whose scope
contains only $I$ still checks $A$ and $B$, and a level computed with the stale
divisor is refused.

\emph{Well-definedness.} \texttt{withSnapshot} (\S\ref{sec:basis-types}) fixes $\beta$
as $\beta_t$ restricted to the \emph{stamp-closure} of $S$ --- the set
$S\cup\bigcup\{\mathrm{dom}(\mathrm{stamp}(d)) : d \text{ consumed}\}$. This restriction
is well-defined, because $\beta_t$ is total over registered units and \texttt{ingest}
admits no stamp naming an unregistered unit. Thus $S$ restricts which balances are read,
never the domain on which admissibility is checked.

Clause (iii) is what makes the invariant meaningful for a black box: the
framework legislates which data it is legal to feed a pricing function, never
what the function does with them. State-sufficiency (\S\ref{sec:valuation})
presupposes that prices and balances agree in basis, and the snapshot's type
variable carries exactly that agreement. A held unit priced in the wrong basis
remains representable as an error but is unrepresentable as a \texttt{Cash}
value: one more illegal state that the valuation types exclude, beyond those of
\S\ref{sec:valuation}.

The same rule is also recorded as a condition, C13, in the register of
\S\ref{sec:states-conditions}. Its statement keeps the word \emph{arrow} of the formal
register for what this section's prose calls a declared conversion.

\begin{description}[style=nextline]
\item[C13 (basis edge).]\hypertarget{cond:C13}{}\label{cond:C13} Every consumed datum
names its joint basis; consumption requires pointwise agreement, over the stamp's whole
domain, with the prevailing total assignment $\beta_t$, or a ledger-authored arrow. C13 does for a datum's meaning what P3 does for a move's existence: P3 refuses a
move to a wallet that does not exist, and C13 refuses a datum consumed against state
that sits in a different basis.
\end{description}

\subsection{Type-level enforcement}
\label{sec:basis-types}

The guarantee is the two-tier discipline of \hyperref[cond:C11]{C11}: inside
authored code the basis is a type index; in the store it is plain stamped data;
at the read seam the index is re-established by checking the stamps. Which
epoch a unit occupies is known only by replaying the committed log at run time,
so the compiler cannot know an epoch's value. The honest static guarantee is
therefore \emph{sameness} --- two values carry the same index --- never the
index's value. A type-level ordinal in the normative layer would claim static
knowledge the system cannot have.

\paragraph{Tier 1 --- the didactic law.} The listing below states the invariant
on one page; it is not the system mechanism. It gives a price the type
\texttt{PriceAt e} and a balance the type \texttt{BalAt e}, where \texttt{e} is
an epoch index that appears only in the type, never inside the value --- such a
parameter is called a \emph{phantom} parameter. A quantity and a price meet
only at equal index, and the only function that changes an index applies a
ledger-authored adjustment. All constructors are unexported, so the sole author
of an \texttt{Adj} value is the ledger's fold; an \texttt{Adj} value is
therefore itself a proof that the log carries exactly these boundaries, in
effective order. The commented \texttt{\_bad} at the end shows the phantom
doubling of the opening as a rejected program: the post-split balance of 2,000
shares priced with the pre-split quote of 100 is not a mistaken number; it is not a
well-typed program at all.

One language fact must be handled before the listing is sound. GHC assigns
every type parameter a \emph{role}, which governs when two instantiations of
the type may be coerced into one another. For a phantom parameter the inferred
role permits coercion between any two indices --- which would forge exactly the
crossing the types exist to forbid. The declaration \texttt{type role \dots\
nominal} withdraws that permission: coercion becomes derivable only at equal
index. The annotations in the listing are therefore normative, not stylistic.

\begin{lstlisting}[language=Haskell]
{-# LANGUAGE DataKinds, GADTs, KindSignatures, RoleAnnotations #-}
data N = Z | S N

newtype PriceAt (e :: N) = PriceAt Rational   -- constructor NOT exported
newtype BalAt   (e :: N) = BalAt   Integer    -- constructor NOT exported
type role PriceAt nominal   -- GHC would infer PHANTOM, and Data.Coerce.coerce
type role BalAt   nominal   -- (Safe Haskell, no unsafe feature) would then forge
                            -- the crossing at ANY pair of indices. Nominal makes
                            -- the coercion derivable only at equal index.

data Adj (m :: N) (n :: N) where              -- constructors NOT exported:
  AIdA  :: Adj n n                            -- the ledger fold is the only author
  AStep :: Adj m n -> OpSpec -> Adj m ('S n)  -- one boundary, in effective order

adjust :: Adj m n -> PriceAt m -> PriceAt n   -- the ONLY epoch-crossing function

markValue :: BalAt e -> PriceAt e -> Cash     -- Qty and Price meet ONLY at equal index
markValue (BalAt q) (PriceAt p) = ...

-- _bad :: BalAt ('S e) -> PriceAt e -> Cash
-- _bad = markValue
--   TYPE ERROR. The epoch-(n+1) position priced with the epoch-n quote is not a
--   wrong number; it is not a program. And AStep extends only to 'S n, so the
--   mis-ordered composite of the June-1 timeline's 49-vs-50 has no inhabitant.
\end{lstlisting}

\paragraph{Tier 2 --- the normative seam.} \texttt{withSnapshot} is the one
constructor of coherent read scopes, mirroring \texttt{applyTx} as the one
write door. The caller hands it a function of type
\texttt{forall b.\ Snapshot b -> r}; the type variable \texttt{b} is chosen
fresh at each call, and the caller can neither name it nor equate it with any
other --- a variable with this status is called a \emph{skolem}. It denotes a
whole basis assignment, not one unit's epoch, and every accessor is unit-keyed.
The ledger is sealed inside the scope: no \texttt{Ledger} parameter escapes
downstream. The per-unit price accessor is fallible, so quarantine is
representable.

\begin{lstlisting}[language=Haskell]
data Snapshot b   -- ABSTRACT. Constructed only by withSnapshot, which fixes the
                  -- target assignment beta = beta_t on the STAMP-CLOSURE of the
                  -- scope (Invariant B: beta_t is total over registered units;
                  -- Scope restricts which balances are read, never beta's domain),
                  -- reads balances from the SAME sealed ledger, and transports
                  -- every admissible stamped observation to beta along declared
                  -- arrows -- refusing, per unit and kind, where none exists.
type role Snapshot nominal

newtype BalAt'   b = ...        -- constructors NOT exported
newtype PriceAt' b = ...
newtype At b a     = ...
type role BalAt'   nominal
type role PriceAt' nominal
type role At       nominal nominal

withSnapshot :: Ledger -> Timestamp -> Scope
             -> (forall b. Snapshot b -> r) -> r     -- quantifier where the proof lives

snapBal   :: Snapshot b -> WalletId -> UnitId -> BalAt' b
snapPrice :: Snapshot b -> UnitId -> Either BasisError (PriceAt' b)  -- per-unit refusal

markV :: BalAt' b -> PriceAt' b -> Cash                  -- the seam; Cash is basis-free
zipAt :: (x -> y -> z) -> At b x -> At b y -> At b z     -- in-scope derivation

data Valuation = Priced Cash
               | Unpriced UnitId BasisError              -- quarantined, enumerated
               | PricedStale UnitId BasisId BasisId Cash  -- logged election

value :: Snapshot b -> [WalletId] -> (Cash, [Valuation]) -- book total + exclusions

-- The single cross-scope door, ledger-authored, factoring ONLY through declared
-- arrows:
transport :: Snapshot b0 -> Snapshot b1 -> Kind a -> At b0 a
          -> Either BasisError (At b1 a)
\end{lstlisting}

\paragraph{Unreachability, for an arbitrary black-box pricing function.} Let
$F::\forall b.\ \texttt{Snapshot } b\to r$ be any pricing function; nothing is
assumed about its body. Because $F$ must work for every \texttt{b}, and the
fragment stated below closes every side channel, the type alone constrains what
$F$ can do with \texttt{b}-indexed values. This consequence of the quantifier
is called \emph{parametricity}, and it yields the argument in two steps.

First, every \texttt{b}-indexed value that $F$ manipulates comes from its one
\texttt{Snapshot} argument. The carriers' constructors are unexported, so $F$
cannot build a \texttt{BalAt' b} or a \texttt{PriceAt' b} of its own; the only
ways to obtain one are \texttt{snapBal}, \texttt{snapPrice},
\texttt{transport}, and the index-preserving combinators, and each of these
yields values at the snapshot's target assignment $\beta$ by construction of
the door. A quantity and a price that $F$ combines therefore both descend from
the same snapshot, and both sit at the same assignment.

Second, values from two different snapshots cannot be mixed. Each call to
\texttt{withSnapshot} introduces its own skolem, and two skolems never unify: a
price indexed by one snapshot's variable does not typecheck where the other's
is required. The single function that crosses snapshots, \texttt{transport},
factors through the unique declared conversion or refuses with a typed error.

Hence every (quantity, price-like) pair any well-typed $F$ can form lies in a
single basis assignment. In the small world: $F$ may pair 2,000 shares with 50
inside the post-split snapshot, or 1,000 with 100 inside the pre-split one, but
no well-typed $F$ pairs 2,000 with 100 --- the \EUR{200{,}000} phantom is not
a wrong output of $F$; it is not a program. What the types cannot check is the
store boundary, where stamps are erased data. There the runtime complement is
a runtime check inside \texttt{withSnapshot} that the two bases really are equal, which returns a typed
\texttt{BasisError}, never a silent pick. This is the C11 guarantee restated at
the read seam: compile-time exclusion at authorship, typed refusal at the
boundary --- claimed as exactly that and no more.

\paragraph{Language-fragment proviso.} The parametricity argument holds only in
a stated fragment of the language, because GHC offers features that would
defeat it. The central one is coercion. GHC infers the role \emph{phantom} for a type
parameter that occurs in no constructor field. At role phantom, the
\texttt{Coercible} relation holds between all index pairs. So
\texttt{Data.Coerce.coerce} --- which Safe Haskell accepts, with no unsafe feature
involved --- would convert a value at one index into a value at any other, forging
the epoch and skolem crossings past every unexported constructor. Every
basis-indexed carrier therefore declares \texttt{type role \dots\ nominal} on
its basis index; the annotations are normative, not stylistic. The
unreachability claim above is made for exactly this fragment:
\begin{itemize}[nosep]
\item Safe Haskell;
\item nominal roles on every basis-indexed parameter;
\item constructors, field selectors, and eliminators unexported;
\item no \texttt{Generic}, \texttt{Data}, or \texttt{Typeable}-derived path exposing
representation;
\item no \texttt{GeneralizedNewtypeDeriving} on the carriers;
\item no Template Haskell in consuming modules.
\end{itemize}
Within the fragment the free-theorem reading of $\forall b$ --- the
parametricity argument above --- is sound, and \texttt{coerce} contributes no
way to build or convert a basis-indexed value.

\subsection{Time travel and late corporate actions}
\label{sec:basis-time-travel}

Historical valuation needs no new machinery. \texttt{clone\_at($t$)} already rebuilds
every \texttt{UnitStatus} field by replaying the log up to $t$ (\S\ref{sec:lifecycle}),
and \texttt{usBasis} is such a field, so the clone carries the basis that prevailed at
$t$ (P8, extended). Valuing history is \texttt{withSnapshot} over the clone, at the
target assignment $\beta_t$. Raw observations are immutable and stamped, so they never
migrate. A datum stamped after $\beta_t$ can be carried backward only if every
conversion between the two bases is invertible (Definition~\ref{def:opspec}): the
split's \texttt{Scale}~$2$ is invertible, so a post-split print
reaches a pre-split valuation, while a \texttt{Pending} spin-off declares no conversion,
so nothing crosses and the datum is refused into the workflow of
\S\ref{sec:basis-failure}.

Two version axes run through every historical question, and they are independent. The
\emph{basis axis} records the order in which boundaries take economic effect; it is
ledger-authored, and conversions compose along it
(Proposition~\ref{prop:composition-law}). The \emph{knowledge axis} records what the
ledger knew at each moment; it is bitemporal: an ``as known at $t$'' replay uses the
chain as the log held it at $t$, and a ``restated'' replay uses the chain with
corrections applied. A historical question fixes a point on each axis, and the two
choices compose.

A late or retro-effective boundary books at $t_{\mathrm{notify}}$ but takes economic
effect at the earlier $t_{\mathrm{eff}}$. It changes no stored stamp, because stamps
are stable ids. Instead the ledger appends \emph{re-stamp derivation events} for the
observations captured in $[t_{\mathrm{eff}},t_{\mathrm{notify}})$: each keeps its
value, receives the corrected coordinate, and carries lineage to the late transaction.
That transaction's \texttt{SetBasis} is tip-welded
(Principle~\ref{prin:tip-weld}): it re-asserts the standing effective tip, so the live
coordinate never moves backward. The corrected replay, meanwhile, projects the chain with
the boundary at its true effective position mid-chain and recomputes the conversions that
span it. The two replays then diverge exactly as intended: the as-known replay
reproduces the honest mistake, and the corrected replay applies the re-stamps. Moves
committed inside the window repair through the existing compensating-transaction path
(\S\ref{sec:implementation-corrections}), now with a citable coordinate lineage. Pinning both the snapshot and the chain
version gives a valuation that is both
repeatable and correct; pinning the snapshot alone, as the old machinery did, repeated
the earlier mistake faithfully instead of correcting it.

\subsection{Failure semantics as workflow}
\label{sec:basis-failure}

Every refusal carries a typed error, a \texttt{BasisError} value. What the refused
datum was for decides what happens next: the consequence class selects the regime. No
regime is silent.

\renewcommand{\arraystretch}{1.25}
\begin{center}
\begin{longtable}{@{}p{2.2cm} p{3.6cm} p{4.7cm} p{2.9cm}@{}}
\toprule
\textbf{Regime} & \textbf{Trigger} & \textbf{Consequence} & \textbf{Mandatory for} \\
\midrule
\endhead
W1 Block & Move-emitting consumption (settlement projection, NAV strike, fee
crystallisation) meets any \texttt{BasisError}. & The lifecycle event defers as a
recorded workflow event with a named unblock condition --- a fresh observation at
$\beta$, or the crystallising publication event --- queued under obligation liveness
(P21, \S\ref{sec:obligation-liveness}). Real money never moves on an unwitnessed basis.
& The stale-data gate (\S\ref{sec:implementation-faults}) and the pre-invocation check
(\S\ref{sec:temporal}): basis equality first, timestamp threshold second. \\
\addlinespace
W2 Per-unit quarantine & \texttt{snapPrice} returns \texttt{Left} for a unit: no
admissible datum, a \texttt{Pending} gap, an undeclared kind. & The book values on time
as $(\texttt{Cash},[\texttt{Unpriced}\ \dots])$ --- a total minus a named, enumerated
exclusion list. Blocking the whole book is rejected: one late notice must not halt
thousands of units; shippability decides. & Daily marks, risk, management PnL. \\
\addlinespace
W2$'$ Flagged-stale carry & An explicit, logged election event with its canonical
writer row in the C11 table. & The unit is valued \emph{coherent-at-}$\beta^-$ --- the
latest assignment $\beta^-\le\beta_t$ at which its data are pointwise complete --- never
a mixed vector. The result is \texttt{PricedStale}, a distinct constructor carrying its
gap; downstream must pattern-match. \emph{Coherent-at-}$\beta$ means all inputs share
assignment $\beta$; \emph{current} means $\beta=\beta_t$; W2$'$ is coherent, not
current, and says so. & Management PnL where quarantining a large unit is worse than a
flagged carry. \\
\addlinespace
W3 Partition quarantine & Aggregation partitions sources by asserted basis id;
quorum and divergence are computed within a partition only. & $\{100\ @\ b_3\}$ and
$\{50\ @\ b_4\}$ are two singleton partitions --- no false divergence, no median 75 that
exists in no basis. A large innovation with no committed boundary is the signature of an
undeclared corporate action or a bad feed: the series quarantines. This is TA-BASIS's
detection signal; the window in which the world has split but the log has not is
detected, not prevented. & All ingestion around corporate-action windows. \\
\addlinespace
W4 Notice attestation & A boundary declaration rewrites both sides of the seam and,
once its moves are emitted, cannot be repaired by the price-correction path. & Admission
requires issuer or exchange signature, or content-hash agreement across
$N_{\mathrm{CA}}\ge2$ independent sources (owner: data governance). Same-$t_{\mathrm{eff}}$
weld: a boundary whose $t_{\mathrm{eff}}$ equals that of another boundary on the same
unit --- committed or co-presented --- is admitted only if the attested notices declare
mutual precedence ($\mathit{prec}$, Definition~\ref{def:basis-id}); absent
declaration, the colliding notice is refused and the unit enters pending-transition,
fail-closed. Between
first notice and confirmation the unit is \emph{pending-transition} for
consumption across the boundary: W1 blocks, W2
quarantines, W3 partitions. & Every boundary event. \\
\bottomrule
\end{longtable}
\end{center}

\subsection{The 2-for-1 split, end to end}
\label{sec:basis-worked}

A wallet holds 1{,}000 shares of $u$, and $\texttt{usBasis}=b_3$.

\begin{enumerate}[nosep]
\item \emph{Snap.} Vendor EX1 publishes 100.00. The ingest door stamps the observation
$\{u\mapsto b_3\}$, per the standing convention for the pair $(u,\mathrm{EX1})$. The
stored record is immutable and logged.
\item \emph{The corporate action.} One atomic transaction $\tau$ carries three things:
a move of $+1{,}000$ shares from the corporate-action virtual wallet
($1{,}000\to2{,}000$), the status write \texttt{SetBasis}\,$b_4$, and the declaration
$\{\mathrm{KSpot}\mapsto\texttt{Scale}\ 2,\dots\}$. Admission checks the welds:
the quantity legs satisfy $2{,}000=1{,}000\times2$ and the cash legs satisfy
$0=1{,}000\times0$ (Theorem~\ref{thm:basis-value-invariance}), and $b_4$ is the
effective tip after insertion (Principle~\ref{prin:tip-weld}). No reachable state
holds the doubled quantity without the boundary that explains it.
\item \emph{Snapshot at $t_{\mathrm{price}}$.} The assignment gives $\beta(u)=b_4$; the
fold gives a balance of 2{,}000 at $b_4$; the stored spot is stamped $b_3$. The door
carries the spot across the one declared conversion: $A(b_3\to b_4)(100)=50$.
\item \emph{Mark.} $\texttt{markV}\ 2{,}000\ 50=\EUR{100{,}000}$. No phantom.
\item \emph{The bug, attempted.} Feeding the raw 100 into the mark requires a value of
type $\texttt{PriceAt'}\ b$. The only functions that produce one are \texttt{snapPrice}
and \texttt{transport}, and each either refuses or converts. The basis index
carries a nominal role annotation (\S\ref{sec:basis-types}), so the type is
distinct at distinct bases and a raw number cannot be coerced across the
boundary. In the didactic Tier~1 listing of \S\ref{sec:basis-types}, pairing
the epoch-$(n{+}1)$ balance with the epoch-$n$ quote does not compile. The \EUR{200{,}000}
is therefore a program that does not exist on the compile side and a
\texttt{Left BasisMismatch} on the runtime side; it never reaches a committed
\texttt{Cash}.
\item \emph{Attribution.} The true post-split spot is 55 --- up 10\% like-for-like.
Open one \texttt{Snapshot} pinned before the split and one pinned after it;
\texttt{transport} carries the opening pair along the declared conversion,
$(1{,}000,100)\mapsto(2{,}000,50)$, and Theorem~\ref{thm:basis-value-invariance} with
$c=0$ makes the carried value equal. Then
$\mathrm{PnL}_{\text{price}}=2{,}000\times(55-50)=+10{,}000$ and
$\mathrm{PnL}_{\text{flow}}=(2{,}000-2{,}000)\times55=0$. A naive decomposition that
compares the \EUR{55} close against the un-transported \EUR{100} open instead splits
the \EUR{10{,}000} into $1{,}000\times(55-100)=-45{,}000$ of ``price'' and
$(2{,}000-1{,}000)\times55=+55{,}000$ of ``flow'' --- two large terms that are artefacts
of mixing the pre- and post-split frames. Transporting the opening pair across the split
removes them: the decomposition, not only its sum, is basis-invariant, and it can be
written only through the one declared door.
\item \emph{Late notice.} If $\tau$ is not yet logged at $t_{\mathrm{price}}$,
\hyperref[cond:C3]{C3} keeps the moves and the boundary together, so the snapshot is
coherent at $b_3$: $1{,}000\times100=\EUR{100{,}000}$ --- the right answer under the
information held. Meanwhile the exchange disseminates prices near 50: W3 quarantines
the series, and W1 blocks settlement-grade use. At $t_{\mathrm{notify}}$, $\tau$
commits with economic time $t_{\mathrm{eff}}$. If boundaries with later
effective times committed in the interim, four mechanisms reconcile them, each
already stated. The tip weld (Principle~\ref{prin:tip-weld}) makes $\tau$'s
\texttt{SetBasis} re-assert the standing tip, so the coordinate never moves
backward. The chain projection places $\tau$'s boundary at its effective
position, mid-chain. Re-stamp events (\S\ref{sec:basis-time-travel}) move the
window's observations to their corrected coordinates. Replay of the corrected
log is then coherent end to end. At no point does $2{,}000\times100$ or
$1{,}000\times50$ reach a committed \texttt{Cash}.
\end{enumerate}

\noindent The same split read from the boundary --- two consumers snapping one series
on opposite sides of the open --- is the worked example of the market-data contract
(\S\ref{sec:basis-contract}).

\paragraph{The dividend.} The same walk on a 1{,}000-share holding priced at 102, for
a \EUR{2} dividend: the declaration is \texttt{Shift}~$-2$ with quantity factor $f=1$.
Transport gives $102-2=100$. The 2-per-share cash leg books once, in the same
transaction, and the weld checks $\text{booked cash}=q\cdot c$:
$1{,}000\times2=\EUR{2{,}000}$, so the total, \EUR{102{,}000}, is unchanged.

This recovers the \texttt{Cum}/\texttt{Ex} pricing of
\S\ref{app:pricing-coordination} exactly: that appendix treats this very
dividend, with $f=1$ and a chain only one boundary long. Its quote-ex flag
and this section's per-(unit, source) stamp carry the same one bit of
information --- whether a quote is cum or ex --- read from the two sides of
the door. In the cum-to-ex direction, a source still quoting cum past the
ex-date is the lagging convention of \S\ref{sec:basis-ingest}, and the seam
converts its prints by $p-c$ on read. In the ex-to-cum direction, a source
quoting ex before the ledger's transition is the one case the appendix's flag
cannot express. Here it is a large innovation --- a price jump no committed
boundary explains --- and W3 quarantines the series
(\S\ref{sec:basis-failure}), fail-closed.

The bond coupon is the same event with $f=1$. The coupon date books the cash
move and advances the series from cum to ex in one atomic transaction
(\hyperref[cond:C3]{C3}). Decoupling the two admits exactly two corruptions.
Cash paid while the mark stays cum counts the coupon twice; the mark gone ex
with no cash booked loses it. Atomicity makes both unreachable, because the
price jump at the coupon date exactly offsets the booked cash
(Theorem~\ref{thm:basis-value-invariance}).

\paragraph{The index divisor.} Index $I$ is computed over constituents A and B. With A
at 60 and B at 84, the published level $(60+84)/D=100$ fixes the divisor $D=1.44$,
snapped as a derived observation whose joint stamp covers everything it consumed:
$\{I\mapsto b_I,\ A\mapsto b_A,\ B\mapsto b_B\}$. B then splits 2-for-1
($b_B\to b_B'$), and fresh prices arrive: A 60, B 42. The corrupt level
$(60+42)/1.44\approx70.83$ --- a phantom $-29.2\%$ index move --- can be computed only by
consuming the snapped $D$ at an assignment where $\beta_t(B)=b_B'\ne b_B$. That
consumption is pointwise inadmissible --- the stamp must agree with the
assignment at every unit it names, and it disagrees at B --- and is refused
fail-closed. The check is well
defined even though the valuation scope holds only $I$: $\beta_t$ is total over
registered units, and admissibility ranges over the stamp's full domain $\{I,A,B\}$
(Invariant~\ref{inv:basis}(ii)). The framework will not transport a derived datum
(Principle~\ref{prin:rederivation}). Three legal paths remain: the index authority's
maintenance boundary on $I$ declares \texttt{Recompose} with $D'=1.02$, so
$(60+42)/1.02=100$, entering under W4 like any corporate action; or a fresh divisor
observation arrives stamped at the new assignment; or the level is re-derived in scope
from transported constituents. Until one of these happens, $\texttt{snapPrice}\ I$
returns a \texttt{Left}: \texttt{Unpriced}~$I$, enumerated, never a silent 70.83.

\paragraph{The delisting.} A wallet holds 1{,}000 shares of $u$; the last
stamped print is \EUR{12.00}. The confirmed notice declares liquidation with
no successor and no consideration: class \texttt{CATermination}, declaration
\texttt{Terminal} on every kind, any residual distribution booked as terminal
moves in the same transaction. \texttt{Terminal} gives no conversion in
either direction (Definition~\ref{def:basis-category}), so the stale 12.00
cannot cross: \texttt{snapPrice}~$u$ returns a \texttt{Left}, and the book
values as $(\texttt{Cash},[\texttt{Unpriced}\ u])$ (W2). The \EUR{12{,}000}
stale carry never reaches a settlement-grade use. That is not a failure to
price; it is the only correct mark for a dead series, since any number
produced from the stale print would be a guess. The flagged stale carry
(W2$'$) remains available for management PnL, and it arrives as
\texttt{PricedStale}, pattern-matched downstream --- never a silent 12.00.
Termination is a fact of the confirmed notice, never of the venue. A series
that leaves its exchange but continues over the counter has no boundary event
at all, because no per-share economics changed: what ends is the exchange
source convention. The feed goes silent, the stale-data gate and W2 do their
work, and fresh prints re-price the same unit at the same basis under a new
(unit, source) convention (O4).

\paragraph{The spin-off.} A wallet holds 1{,}000 shares of A, stamped
\EUR{100} at basis $b$. The confirmed notice: one share of A1 and one of A2
per share of A, effective at the ex-date; class \texttt{CASpinOff}. When the
children later print on their own, A1 will open near \EUR{80} and A2 near
\EUR{20}. Neither number is in the notice, because no notice can contain a
price formed after it. The confirmed terms therefore carry exactly as far as
they logically can, and the two planes carry different halves.

On the entitlement plane, moves issue 1{,}000 A1 and 1{,}000 A2, the children
registered at their origins with their datum kinds (\S\ref{sec:basis-kinds}).
On the quotation plane, the declaration is a \texttt{Recompose} of old A onto
the basket $\{A1\mapsto1,\,A2\mapsto1\}$. The basket mechanics, pointed to
from Definition~\ref{def:opspec}, are these. A \texttt{Recompose}~$\sigma$
whose target names successor units declares one forward conversion. It
carries the parent's price-like kinds onto a basket datum, jointly stamped at
the successors' origins; the basket weights are the confirmed entitlement
ratios. The basket datum is consumable only jointly: pointwise stamp agreement
(\hyperref[cond:C13]{C13}) refuses its projection onto any single component.
No backward or per-component conversion exists until fresh per-component
stamped observations arrive, because the decomposition is not in the notice.
The transported basket value carries its lineage (\texttt{txCause}, the
declaring boundary). And the basket step is a bridge, not a standing
conversion: it retires on arrival of fresh per-component observations, after
which composition proceeds component-wise under
Proposition~\ref{prop:composition-law}.

Pair-level valuation is therefore all the stale \EUR{100} mark can license;
component pricing stays refused. The stale \EUR{100} may value the pair --- 2{,}000
shares combined, marked through the basket at
$1{,}000\times100=\EUR{100{,}000}$, phantom-free, so the overnight risk of
the combined position is readable at the open. It may never value a
component: C13 refuses the projection of the joint stamp onto A1 or A2 alone,
because that projection is the 80/20 decomposition the notice never
contained. Per-component kinds stay \texttt{Pending}, and the pending
spin-off is the normal case, not the edge case: the basis id advances at
effectiveness, and until fresh per-component stamps arrive no number can
price a component at the new basis (P34). Once A1 and A2 carry their own
stamps --- the 80 and the 20 --- the basket step has retired, and the book
marks $1{,}000\times80+1{,}000\times20=\EUR{100{,}000}$ from observations,
component-wise on each child's chain.

The refused error, in numbers: carrying the stale 100 onto A1 by an assumed
identity (\EUR{100{,}000}) and adding the fresh A2 print at 20
(\EUR{20{,}000}) books \EUR{120{,}000} against a true \EUR{100{,}000} ---
A2's value counted twice, once embedded in the stale parent price and once as
the fresh print. The seam refuses it twice over. \texttt{AId} is declared,
never assumed, so no conversion carries A onto A1 alone. And C13 refuses the
mixed consumption of a stale pair value beside a fresh component print.

\paragraph{The elective merger.} A wallet holds 1{,}000 shares of target T,
stamped \EUR{100}; the acquirer is Q. The confirmed terms after election and
proration: per T share, \EUR{102} cash or 2.04 Q shares, under a 50\%
aggregate cash cap. Holder A (500 shares) elected cash, holder B (500 shares)
elected stock, and the cap does not bind.

Resolution comes first. All of it is in the log, and none of it is an
operator. The announcement opens the window, and the unit enters
pending-transition, so consumption across the boundary fails closed.
Election instructions are logged per holder. Proration is computed.
Confirmation publishes the final blend. Only now does an operator exist.

On the entitlement plane the two holders receive different moves: A's 500 T
become \EUR{51{,}000} cash, B's 500 T become 1{,}020 Q. On the quotation
plane the series has one fate. The confirmed blend per T share is \EUR{51}
cash plus 1.02 Q, and the declaration on T's price-like kinds is
$\texttt{Subst}(Q,1.02)\circ\texttt{Shift}(-51)$ --- the composite the
benchmark-cessation fallback of \S\ref{sec:basis-kinds} already uses, no new
case in the menu. Transport of the stale datum:
$100\mapsto(100-51)/1.02=2450/51$ per Q share, exact rational.

Admission is the elective case of Principle~\ref{prin:invariance-weld},
stated here, where its data are on the page. \emph{Elective aggregate weld}:
the confirmed notice declares an election structure and a proration result
(\texttt{CAMergerElective}), and the transaction is admitted only if the legs
in aggregate satisfy $\sum(\text{quantity legs})=f\cdot\sum q$ and
$\sum(\text{cash legs})=c\cdot\sum q$ at the confirmed blend $(f,c)$, up to
explicit cash-in-lieu legs. Here the check reads: quantity legs
$1{,}020=1{,}000\times1.02$; cash legs $\EUR{51{,}000}=1{,}000\times51$. No
single holder's legs match either factor --- that is what an election means
--- and the blend exists exactly because proration is confirmed.

Value before, $1{,}000\times100=\EUR{100{,}000}$. Value after,
$1{,}020\times2450/51+51{,}000=49{,}000+51{,}000=\EUR{100{,}000}$ ---
continuous to the minor unit. Per holder the values differ by the elected
consideration: A holds \EUR{51{,}000} against \EUR{50{,}000} before, and B
holds \EUR{49{,}000}. The \EUR{1{,}000} A gains is the \EUR{1{,}000} B cedes,
determined by the logged elections and summing to zero. If Q's next fresh
print is 50, the transported $2450/51\approx48.04$ is not wrong: it is old
information in new coordinates. The gap --- \EUR{2{,}000} on the
\EUR{100{,}000} aggregate --- is marking PnL, old information against a fresh
mark, not a phantom. The operator never absorbs it.

\subsection{CDM corporate-action mapping}
\label{sec:basis-cdm}

Real corporate actions arrive as a message pipeline: announcement, election
instructions, confirmation. The pipeline is not mechanical. Terms are
indicative at announcement. Per-holder elections and proration change what is
delivered. Mixed consideration blends cash and stock. All of that is lifecycle
\emph{resolution}: events in the log, entitlements as moves, no operator. The
message standards that carry the pipeline --- ISO~15022 (MT564/565/566),
ISO~20022 (seev.031/033/036), or a vendor feed --- are sources under the W4
attestation regime (\S\ref{sec:basis-failure}). They are named here as
instances, never as a type layer: each standard is one more attested source
under rules that already exist, and no message format adds a type or a rule
to this specification. ``Confirmed'' is a property of attestation, not of
format: a notice is admitted on its W4 attestation alone, whichever format
carried it.

An operator exists only from confirmation. This is the resolution/consumption
boundary, hereafter the \emph{confirmation gate}: ``the declared terms of the
notice'' in Definition~\ref{def:opspec} means the \emph{confirmed} terms ---
the terms carried by the W4-attested final-terms notice. The attested object
is terms-confirmation; movement confirmation --- the message class that
reports postings after the fact --- has no gating role, since it can post days
after the ex-date and gating on it would quarantine every split from ex-date
to pay-date. The boundary's $t_{\mathrm{eff}}$ is the ex-date. Between
announcement and confirmation no operator has been declared, so nothing
adjusts: consumption that would cross the announced boundary fails closed or
quarantines, while consumption in the unit's standing basis --- fresh prints
before $t_{\mathrm{eff}}$ --- is unaffected, since the basis advances only at
effectiveness. The unit is pending-transition for the crossing (W4,
\S\ref{sec:basis-failure}).

The sentence of
\S\ref{sec:basis-operators} --- an event that declares no operator adjusts
nothing --- is exactly this window's safety clause, and P32 witnesses it:
pending blocks consumption, and an empty declaration yields refusal, never
identity.
Confirmation is where the declaration of
Principle~\ref{prin:one-decl-two-uses} is written; both of its uses begin
there. For mandatory actions --- splits, ordinary dividends --- final terms
confirm before the ex-date in the ordinary course, so the window is empty and
consumption never blocks; the window bites exactly where marking would be
guessing: elective and prorated events, late notices, contested terms.

The pipeline's guarantees restate in four steps. The operator is
\emph{declared}: $F$ --- the map of \S\ref{sec:cdm} that turns a business
event into a ledger transaction --- derives it from the notice's confirmed
terms, never from its name (Definition~\ref{def:opspec}). The terms are
\emph{attested}: they are the W4-attested final terms, and no other terms
count. The operator is \emph{welded}: admission checks it against the booked
moves (Principle~\ref{prin:invariance-weld}). And it is \emph{recomputable}:
any party recomputes it from the logged terms and diffs (P27). Declared,
attested, welded, recomputable.

Every confirmed notice declares its class, drawn from a closed taxonomy of
eight:

\begin{lstlisting}[language=Haskell]
data CAClass
  = CASplit             -- split, reverse split, stock dividend: any pure ratio event
  | CADividendOrdinary  -- ordinary cash distribution, coupon detachment
  | CADividendSpecial   -- special or extraordinary distribution, return of capital
  | CASpinOff           -- distribution of a new series to holders of the parent:
                        --   spin-off, demerger, rights issue (the right is the registered child)
  | CAMergerStock       -- stock-for-stock succession
  | CAMergerElective    -- elective or mixed consideration, tender with proration
  | CATermination       -- merger for cash, liquidation, delisting with termination
  | CARebalance         -- index composition or divisor change
  deriving (Eq, Ord, Show, Enum, Bounded)
\end{lstlisting}

\noindent
The class is declared data on the attested notice --- the issuer's or
exchange's classification --- never inferred from a name; the operator still
comes only from the declared terms (Definition~\ref{def:opspec}). Ordinary
and special dividends are two classes because the adjustment-schedule rule
language reads only the declared parameters: an ordinary \EUR{2} and a
special \EUR{2} dividend are indistinguishable by $(f,c)$, yet listed
conventions treat them oppositely, and only the class lets a schedule say so
(Principle~\ref{prin:schedule-authority}). \texttt{CATermination} unifies
three names because their fate is identical: \texttt{Terminal}. The
consideration amount is notice data, not a class.

A rights issue is \texttt{CASpinOff}: the right is a tradable series
distributed pro rata to holders of the parent. The right registers as the
child and is priced only by fresh stamps, never by a theoretical value. Its
later exercise is subscription activity of the right itself --- a lifecycle
of that unit, not a corporate action on the parent. An issuer-level scrip
election is the elective class, whose content is any declared election
structure with a confirmed proration blend. A pure stock dividend is
\texttt{CASplit}. A reinvestment plan is a standing instruction executing as
market purchases after the ordinary distribution, and declares nothing.

Two events are excluded
on purpose: delisting with continuation, a venue fact with no boundary event
(\S\ref{sec:basis-worked}); and non-corporate-action terms amendment, the
table's last row, which declares nothing. Extending the taxonomy is a change
to this specification, never a change to data: a notice outside the eight
classes is refused at the door, fail-closed. The schedule totality gate
\hyperref[cond:C14]{C14} quantifies over exactly this closed set, and each
row of the table below carries its class.

One wall stands over the whole pipeline. An operator re-expresses a datum ---
the same fact in new coordinates --- and never updates its information
content. The repricing an action itself causes is new information: implied
volatility collapses on a takeover announcement, the target pins to the
offer, a spun child finds its level. New information enters only as fresh
stamped observations through the ingest door
(Principle~\ref{prin:stamping-authority}); under the model-free directive
(\S\ref{sec:basis-operators}) the discipline never manufactures it by
transforming old observations. Transported data mark the book at old information in
new coordinates; the gap to the next fresh print is market PnL and appears
as such, never absorbed into the boundary.

CDM \texttt{BusinessEvent}s already map to transactions through the map $F$ of
\S\ref{sec:cdm}, which keeps a business event's economic content and drops its
CDM-specific wrapping; ``CDM'' denotes Common Domain Model v6.0.0 throughout, as pinned in
Appendix~\ref{app:cdm-walkthrough}. The extension is one rule. A corporate-action
event's instructions supply both halves of $\tau$: the moves, as today, and the
declaration, keyed on the event's declared terms and never on its taxonomy name. $F$
lands in (moves, \texttt{SetBasis} + declaration) instead of moves alone. It remains
total and deterministic, and the one economic fact it now keeps --- the basis change
the action causes --- is exactly the fact the old map used to discard.

\renewcommand{\arraystretch}{1.25}
\begin{center}
\begin{longtable}{@{}p{4.4cm} p{4.2cm} p{4.8cm}@{}}
\toprule
\textbf{CDM representation} & \textbf{Moves (existing)} & \textbf{Declaration (new)} \\
\midrule
\endhead
StockSplit / ReverseSplit / StockDividend (QuantityChange, ratio $r$) --- class
\texttt{CASplit} & entitlement
$\times r$ via the corporate-action virtual wallet (\S\ref{sec:contracts}) &
$\mathrm{KSpot}\mapsto\texttt{Scale}\,r$, the quantity factor, per-share reference
kinds likewise; weld checks $q'=q\cdot r$ \\
\addlinespace
CashDividend / coupon detachment (Transfer) --- class \texttt{CADividendOrdinary}
or \texttt{CADividendSpecial}, as the notice declares & cash issuer $\to$ holders &
$\mathrm{KSpot}\mapsto\texttt{Shift}\,(-d)$ on cum-quoted kinds; weld checks cash
$=q\cdot d$; subsumes \S\ref{app:pricing-coordination} \\
\addlinespace
SpinOff --- class \texttt{CASpinOff} & issuance moves of the child; child registered
at its origin & parent kinds
\texttt{Pending}\,\textit{ref} until allocation publishes; the publication event
crystallises the parameter; pair-level bridge per the basket mechanics of
\S\ref{sec:basis-worked} (Definition~\ref{def:opspec}), which exits on fresh
per-component stamps (when-issued prints in the ordinary course), not only on
the allocation publication \\
\addlinespace
Merger, stock-for-stock (TermsChange + succession) --- class \texttt{CAMergerStock}
& C8 Breaking amendment + paired
issuance; \texttt{usSupersededBy} & \texttt{Subst}\,$u'\,r$ --- the cross-unit conversion;
predecessor data reach the successor's basis through it \\
\addlinespace
Merger for cash --- class \texttt{CATermination} & terminal moves & \texttt{Terminal}
--- no conversion in either direction;
fail-closed forbids any further crossing \\
\addlinespace
Index rebalance / divisor change --- class \texttt{CARebalance} & rebalance moves on
the strategy unit & boundary on
the index unit via the basis-notice handler;
$\mathrm{KComposition}\mapsto\texttt{Recompose}\,\sigma$, under W4 \\
\addlinespace
Election / proration windows --- class \texttt{CAMergerElective} & per-holder
entitlement differences are moves & boundary
\texttt{Pending} during the window; the unit-level operator crystallises at proration
--- Status stays read identically by every holder
(Principle~\ref{prin:two-planes}) \\
\addlinespace
TermsChange (non-corporate-action amendment) & C6/C8 as today & no declaration; the
basis advances only on declared basis change \\
\bottomrule
\end{longtable}
\end{center}

\subsection*{Key Invariants and Consequences}

\begin{itemize}[nosep]
\item \textbf{One coordinate, one seam (C13).} The corporate-action basis is a field
of \texttt{UnitStatus}, folded from logged boundary events. Every consumed datum
carries a joint basis stamp. Consumption requires pointwise agreement with the total
assignment $\beta_t$, or a ledger-authored conversion. C13 is the data-plane
counterpart of P3.
(\S\ref{sec:basis-coordinate}, \S\ref{sec:basis-invariant}.)
\item \textbf{Value invariance is a theorem, not a property claim.} The invariance
weld makes the identity $(qf)\cdot\frac{p-c}{f}=qp-qc$ a condition of admission, so
marked value is continuous through every boundary up to exactly the booked cash
(Theorem~\ref{thm:basis-value-invariance}) --- per position across uniform
boundaries, in aggregate at the confirmed blend across elective ones. The lifecycle
value-invariance property of
\S\ref{sec:intro} is thereby proved on both factors for basis boundaries.
\item \textbf{No regression, no reordering.} \texttt{SetBasis} is absolute and
idempotent, so P6 is unchanged. The tip weld makes a basis regressed by a late notice
unrepresentable, and replay in booking order agrees with the effective-order chain
projection on \texttt{usBasis} at every prefix of the log
(Principle~\ref{prin:tip-weld}).
\item \textbf{Fail-closed by expressibility.} A declaration is a partial map over
datum kinds: a kind it does not mention has no conversion and cannot cross the
boundary. \texttt{Pending} and \texttt{Terminal} generate no conversion. Identity is
declared, never assumed. (\S\ref{sec:basis-operators}.)
\item \textbf{The phantom is unrepresentable in a well-typed program.} Within the
stated language fragment, every (quantity, price-like) pair that a well-typed pricing
function can form lies in one basis assignment; at the erased boundary the same
discipline is a typed refusal, never a silent pick. A mixed-basis value is
unrepresentable as a \texttt{Cash} and representable only as an enumerated
\texttt{Valuation} exclusion.
(\S\ref{sec:basis-types}, \S\ref{sec:basis-failure}.)
\item \textbf{Two flows at the market-data boundary.} The basis projection flows out:
read-only, versioned, and free of arithmetic. Raw observations flow in through the
single ingest door, where the ledger attaches the stamp. Door soundness (P25) admits
nothing whose stamped basis the log has not committed, and the stamp is a pure
function of pinned recorded data, so real-time, batch, and multi-vendor operation
coincide and no corporate-action logic can exist outside the ledger.
(\S\ref{sec:basis-contract}.)
\item \textbf{Time travel needs no new machinery.} \texttt{clone\_at($t$)} rebuilds
\texttt{usBasis} by the existing replay fold (P8). The basis axis is effective order;
the knowledge axis is bitemporal. Late notices repair through re-stamp events with
citable lineage, and the live coordinate never moves backward.
(\S\ref{sec:basis-time-travel}.)
\item \textbf{Protected elements untouched.} The atomic move, conservation, and log
immutability (P1, P2, P4) hold in unchanged obligation. The whole discipline consists
of a logged event class, the \texttt{SetBasis} status write, admission conditions at
the single door, a named projection, and a typed seam.
\end{itemize}

\noindent
The invariant checks the two numbers together, not one at a time. A quantity and a
price can each be individually correct and still belong to different bases; the phantom
PnL is exactly what multiplying them in that state produces. When both are typed in the
same basis assignment, that product cannot be formed at all, so the phantom PnL is
unrepresentable.
